\quellseite{571}
\chapter[PRACTICAL EXAMPLES]{PRACTICAL EXAMPLES\footnote{Where the name of the writer is not given, the examples are written by the author of this book.—We purposely give not only models after which to pattern, but also some which are designed to evoke criticism.}}

\section{1. Instruction in Biblical History}

\subsection{1. PREPARATION:} \emph{(a) Aim: We would learn today how God made the world and all that is in it.}

Children, how beautiful is this world in which we live! Which season of the year do you like best? Yes, springtime is the most beautiful. Why do we like this season best? Certainly, because everything again becomes fresh and green. The snow melts, the grass grows in the meadows, the trees shoot out their leaves, the flowers blossom, and the birds return again with their songs. How did nature appear in winter? Yes, and the snow was like the white cover which is placed over a dead person. Then what happened to all nature that seemed to be dead? It was as though born again, made new all over. What do we hate to do when everything outside looks so fresh and inviting? The house seems too small, and we eagerly hasten out of doors, breathe deeply the fresh air, and try to outdo the birds in their joyful song. Who makes everything so new and fresh and beautiful in springtime? Correct, Frank; we could never have caused snow and ice to disappear. But the almighty Lord has done this. He let the sun shine more warmly, which caused the ice to break, the snow to melt, the grass to grow, the flowers to bloom, and all things to become like new. God does this every year, but there was once a time when He did much more. He created this great, beautiful world out of nothing. Think of it! What a wonderful deed! We could not cause the smallest blade of grass to grow or even make the tiniest grain of sand, but God made heaven and earth—this entire world, out of nothing. We would learn about this today—how God made heaven and earth and everything that is in it.

\emph{(b) Aim: How the Saviour once helped His disciples on Lake Gennesaret.}

Who of you has ever seen a big lake? Well, you have certainly seen a small lake or a pond. The Lake of Gennesaret is eleven or
\quellseite{572}
twelve miles long (that is about as far as from here to . . . .), and about six miles wide (about as far as . . .? That’s a lot of water, is it not? And it is very deep; and it has a wonderful deep blue color. If one wished to go from one shore to the other, one had to take a boat. Who of you has ever had a boat-ride? Tell the others about it. At the time of which our lesson speaks there were no steamships or gasoline boats, but only sailing vessels. A large cloth was raised on a mast (i. e., a large pole), and the wind blowing against the cloth moved the boat. Picture of a boat on the Lake of Gennesaret. Today we shall hear how Jesus was once on the Lake of Gennesaret, and helped His disciples out of a very great trouble.—Fankhauser, page 55.

\emph{(c) Aim: How the Lord Jesus was brought to a child sick unto death.}

A child “sick unto death”—what does that mean? You have seen a sick child. But have you ever seen a child “sick unto death”? It is a child so sick that he must die. Is that not terrible? I have known children who were so sick that they could not get well, but had to die. Have you ever seen or heard of such? What was the matter? “One had scarlet fever.” What did the poor parents do? (Sent for the doctor.) What did the doctor do? (Watched the child, examined him, and wrote a prescription.) Why did he do all that? Would that make the child better? (Sometimes—but not always.) The poor parents! Why do I call the parents poor? (They are in great trouble and very sad.) What does the mother do? (She sits at the child’s bedside.) And when the father comes home, he, too, is worried, and softly asks mother whether the child feels better. Then father and mother look sadly at each other; and as the child grows worse, they think more of the only one who can help when even the doctor’s skill is helpless. Who is that? And what will father and mother then do? (They will fold their hands and pray to the dear Lord.) “Let our dear child live, make him well”—so will they pray. They want to bring God, Jesus Christ the Saviour to their child that is sick unto death. Today we would hear how a father brought the Lord Jesus to the child which was at death’s door.—H. Spanauth, \emph{Praeparationen fuer den Religionsunterricht. Unterstufe}. ²1910 (according to the method of “construction”).

\emph{(d) Aim: How Jesus raised His friend Lazarus from the dead.}

Two weeks ago a member of our congregation died, whom all of you knew. Who was it? Body and soul were separated from each other, and that is what we call death. Where do the souls of pious
\quellseite{573}
persons go? Where is the body placed? Yes, it is placed in a coffin and then taken to the church where the person when alive so often worshipped; after that it is taken to the cemetery. What is then done? Oh, that is a small, narrow room! It makes no difference how rich or how poor, how high or how low, all—all must go to the grave. Death is like a heartless reaper who with his keen sickle mows down everything in his way. He lays all into the grave. How do we make graves? In eastern lands caves were used as burial places; sometimes these were natural caves, formed by nature, but, more often, they were cut out of the solid rock. The dead were laid into these underground rooms. Grandfather, father and son, besides many others were buried there side by side. Who once upon a time bought such a family tomb? Who was the first one whom Abraham buried there? Who was the next? And the next, etc.? We learned a short while ago of another man who bought such a tomb for himself and his family. Whom do I mean? Who was the first to be laid in this cave? Among us the dead are laid in the earth, and covered with earth; that is why we use a casket with a cover or lid. What is done with the lid before the casket is lowered into the grave? Of course, we have another protection for the casket, which likewise has a lid. Why? We do not wish to bury a human being like a beast, which one simply throws into a hole and covers with soil. We do everything possible to keep the body as long as we can, especially if the departed was near and dear to us. In the East no coffins were used. The corpse was wrapped in a large white cloth with sweet smelling spices, in order that any odor from the corpse might be killed; the body was placed upon a bier, or a stretcher, and so carried to the tomb. What was done in order that wild beasts might not find their way to the body? Where is it plainly stated that a great stone was rolled in front of the cave? We also do something that no animals get into the place where our dead lie. What is it? We put a fence around our cemeteries that our departed ones may rest in peace, undisturbed. What should you therefore not do when you go into the cemetery? It is a place of rest and peace, which should not be disturbed by loud talking or noise. When you go to the cemetery, then think of your grandfather or grandmother, of father or mother, or of others who lie there. Their bodies sleep beneath the sod and wait for the great resurrection day. What do we believe as Christians? Do we know when that shall happen? But this we know, that Christ will some time come again and all the dead shall arise. Then our graveyards will be as full of life as they now are full of the stillness of death. Can the Lord Jesus really call the dead to life? He can do so because He is almighty.

\quellseite{574}
We would hear today how Jesus raised His friend Lazarus from the dead.

\emph{(e) Aim: We learn how the most furious persecutor of the Christians is overcome.} (In abbreviated form.)

The Sanhedrin resolved to persecute the Christians. Stephen is accused by Saul and stoned to death. The Christians flee from Jerusalem, for Saul rages against them. Will their life or, at least, their freedom be in danger outside the capital city? Will their persecutors’ follow them into Samaria, Galilee, Syria, Damascus? Saul, the most violent of all their enemies, will seek them there. He thinks especially of Damascus; for that is the place where many of those have gone who believe that Jesus Christ, crucified and risen, is the only way to salvation. But how shall Saul find entrance into Damascus, and who will give him the right to work against the Christians? He has provided for that; for he has letters to the authorities of Damascus which tell who he is, whence he comes, what he wants, that he comes in the name of Sanhedrin, and that the authorities should help him against the Christians. These letters will open all the doors of Damascus to him. Thus equipped, he goes there. Will he be successful? If so, Christianity shall soon be no more. Or will God hinder and overcome the evil intentions of His enemies? He will; for His name shall be made holy and His kingdom shall come. How shall this be done? We hear how He overcomes Saul, the most furious persecutor of the Christians, on his way to Damascus.

\subsection{2. PRESENTATION}

\subsubsection{(a) The Stilling of the Storm} (narrated in sections for small children).

\emph{1. The Pleasant Voyage.}

The Lord Jesus often sailed with His disciples upon the Sea of Gennesaret, not for pleasure, but so that He could preach to the people on the other side of the sea. (Show map.) Sailing over the waters is far more pleasant than travelling by automobile or railroad. There is no jolting or bumping. Gently and smoothly one glides along. Sometimes the boat rocks gently as a cradle, and one is easily lulled to sleep. The Lord Jesus had worked hard all day. What had He done? Walked much, preached, healed the sick. When you have walked to school and have studied hard there, you are tired when evening comes and easily fall asleep. Jesus was tired, and as the boat went gently gliding through the water and the cool breeze fanned His cheeks, He soon fell asleep. He lay in the back part of the ship with His head upon a pillow. What do you do when you come into the house and see your little brother or your grandmother peacefully sleeping? You walk lightly and close the door
\quellseite{575}
gently, because you love the sleeping one and do not wish to disturb his slumbers. The disciples acted in the same way. They loved the Lord Jesus and wanted Him to have His rest. They were careful and did not talk loudly or laugh, lest they should awaken Him. Thus they sailed on through the golden sunshine. It was a delightful voyage. Now I shall briefly tell you the story just as it stands in the Bible, and you may repeat it to me. Narrate according to Matt. 8: 23. 24; Luke 8: 22. 23.

\emph{2. The Storm and the Danger.}

Now I want to tell you how the Lord Jesus helped His disciples. Do you need help when you are playing? No, but if you have fallen into the river then you need help. When things go wrong help is necessary. But how could matters have gone wrong with the disciples? This morning before Sunday school I showed you the watering trough. The water was smooth as a mirror, but when you blew upon it, it became rough with waves. On the Lake of Gennesaret the wind came suddenly and blew strongly, much worse than if all of you would blow together. Besides, the Lake of Gennesaret is many thousand times larger than a watering trough. The waves rolled high—as high as a house. What happened to the little ship? It was tossed to and fro like the nut-shell in the watering trough. It was no longer pleasant in the boat. Perhaps you have been sitting in a wagon when the horse was frightened and ran away. So the disciples became fearful. It seemed as if their boat must be upset and they must drown. They let the sail down and dipped the water out. They had to hold fast to keep from being thrown into the angry sea. Dark clouds, thunder and lightning, the raging sea,—all threatened death every minute. That was an hour of great need. Narrate according to Matt. 8: 24; Mark 4: 37; Luke 8: 23. Repetition by the children as in first part.

\emph{3. The mighty Helper.}

What do you think? I wonder that the Saviour did not wake up. But He slept peacefully. The disciples saw only the storm and the waves, and in their fear they forgot that Jesus was with them, else they would have thought: He can help us, for He has done many other wonderful things, healed the leper, the man with the dropsy and the like. But they did not think of that. They did not believe that Jesus would certainly see to it that they would not perish. Their faith was weak; they had only a little faith in the Saviour. Finally they awakened Him. What should they have said? They cried: Master, Master, we perish! The Saviour arose and rebuked them, saying: “Why are ye fearful, O ye of little faith!” If you had only believed that so long as I
\quellseite{576}
am with you the ship could not sink, you would not have been so fearful. And now the Saviour helps. How? By dipping out water? By steering to the shore? Oh, no! Listen: “Then He arose and rebuked the winds and the sea, and there was a great calm.” Could any one else have done that? No one, no teacher, no king, no emperor. Had we been there, we should have looked at him in wonder and surprise. “But the men marvelled, saying, What manner of man is this that even the winds and the sea obey Him? Narrate according to Matt. 8:25–27; Mark 4:38–41; Luke 8:24, 25, and repetition by class.—Recital and repetition by class of the entire story in the words of Scripture.—Fankhauser, page 55 ff.

\subsubsection{(b) Zaccheus the Publican}

It was a week before Easter, and the streets of Jericho were filled all the day long with people coming from the North to go to Jerusalem to celebrate the festival of Easter. Naturally, the people of the town did not stay in their houses, but went out on the street; and if they had time, they walked as far as the city gates in order to see the pilgrims and talk to them. In his booth by the city gate sat Zaccheus, the chief tax-collector of the town. During these busy days he had to be at the gate constantly and have his eyes everywhere. He was a small man, well-fed and well-dressed. The tax-collectors under him worked in silence; they did not like him at all. He was a rich man; and he was hated by the people because he made the taxes too high and loaned out money at too high rates of interest. But for some time he had acted as though he cared but little for his fine business and had to force himself to attend to it. He hung his head as if he were really upon the publicans. It was not pleasant when he was about.

Zaccheus also worked in silence, but he was not so much taken up ashamed before the Pharisees and other people who looked down with his accounts as in former times. Again and again he stopped to listen to the talk of the passers-by; and he seemed to be frightened. Why did they all mention one name that caused his heart to stand still? He heard it again: “The Prophet of Nazareth”, “Jesus the Nazarene”. Was He coming to the Easter festival, and did they expect Him at the gates of Jericho? Zaccheus was worried, for he had suffered for quite a while because of that name. A wanderer from Galilee had once come while he was dealing with a poor debtor, and had spoken to him so earnestly and scornfully: Have you never heard what the Prophet of Nazareth teaches: “What shall it profit a man if he gain the whole world and lose his own soul?” That
\quellseite{577}
word had struck deep into his soul, and since that time he had asked more than one traveller about the Prophet, so that by this time he had a good idea of Him and knew much about His teachings and miracles. So it came that, more and more, he felt himself to be hateful and filthy and mean in comparison with Jesus; he had begun to hate his calling, which made him rich through its dishonest dealings. But he could not give up his fine clothing, expensive furniture, soft bed, rich wines and luxurious meals. And besides, what would his wife and children say if they suddenly found themselves without all their luxuries and compelled to live like ordinary working people. It would simply not do:—such had been his difficulties for quite a while. The pencil trembled in his hand:—a little while ago he had ceased writing; his thoughts were far from his work. Then from out the jabbering multitude he heard once more the name of Jesus. He listened intently. Oh, if he might only see Him once—the good, the pure One! Perhaps it would help him! Then he heard a newcomer greeted with the words: “Have you heard, Matthew, that the famous Nazarene is coming?” “What—the Prophet?” Then they whispered, but Zaccheus caught a word that made him tremble. Did they not speak of the Messiah? The mocking words and laughter of another man came to him: “Messiah! Well, I tell you, I expect a different Messiah; my Messiah will wear the finest clothes and command a hundred thousand soldiers! But this Galilean prophet with His few disciples—fishermen and women, yes, even publicans and sinners!” Zaccheus listened more intently. They had said of the Messiah that His followers were publicans and sinners—publicans and sinners like he! Then he must see Him! Breathing deeply, he closed his book and left the gate, while his clerks sneered and scoffed as they watched him depart.

He passed through the crowd as quickly as possible, stopping to listen only when he heard the name of Jesus mentioned. He learned that Jesus was expected toward evening. Reaching his home, he walked restless to and fro. He felt that the greatest hour of his life had come. His wife became worried as she watched him, and finally laid her hand upon his shoulder and asked: “What troubles you?” Then he told her as well as he could what he knew of Jesus, and his desire to see Him. She trembled and folded her hands, and perhaps she imagined that this day might bring a great change into her life. After a little while, Zaccheus went out. He had clothed himself in his finest garments, not for show, but in honor of the Prophet. The
\quellseite{578}
streets were crowded, and with difficulty he made his way to the gate. There the multitude was greater. All Jericho was there, and besides, many of the Jerusalem pilgrims had remained to see the Prophet. So closely were they crowded that one might think a king was coming. Doubtfully Zaccheus looked about. He was hemmed in by the crowd and jostled by many elbows. He reached only to their shoulders. His costly clothing was already torn, and it was impossible for him to force his way through and to see anything. Should it all be in vain? Then he looked up and saw a nut-tree which stood on the side of the street and stretched its branches wide over the crowd. He was only a few steps away from it, but it caused him much effort to reach it. He looked at the slender trunk; he was not used to climbing; but, perhaps, he could make it. Quickly he gathered his garments about him, grasped the trunk of the tree, and began to climb.

The onlookers laughed loudly at the sight of the fat little man with the gold bracelets tearing his shimmering yellow and red clothes on the twigs of the nut-tree. One of them recognized him. “Look at the fat publican Zaccheus,” he cried, and the crowd joined with him in his jeering laughter. With flushed and sweating face Zaccheus reached a place from where he could see. His long hair clung to his face, his clothing was dirty and torn; his white hands scratched and bleeding. The people mocked him, but he cared not, for he had reached his goal and could see the One whom he so much desired to see. From the distance came a murmur which grew to a subdued roar: “He comes! He comes” Zaccheus quivered with excitement. He pushed back his long hair from his forehead, and felt how wet and dirty he was. But what mattered that! There they came—a company of Galileans, easily recognized by their dress. Three men went in advance, and the one in the middle—it must be He! Such eyes full of spirit and power, so much goodness and love in the lines about the mouth—there could not be another such in the whole wide world. He walked so modestly and humbly; He did not look about Him; and yet He moved like a king; every one had to think: If He should command, all would hasten to serve Him. They had drawn near. With a smile one of the men with Jesus touched His arm and pointed out the strange little man who looked down upon them with burning gaze. Those about Him laughed, and a fresh little boy cried: “Yes, yes, that is the rich publican Zaccheus.” Zaccheus flushed hot with shame, and his eyes grew wet with tears. Jesus looked up at him and stopped. He stood there and looked straight into the eyes of the
\quellseite{579}
trembling man. Then He said in friendly tones: “Come down quickly, Zaccheus, for I would sup with you.” Then He turned and went on. All regarded Zaccheus, but none of them laughed. They stared with wonder, envy, astonishment and anger, and many murmured: “So, He will eat with that sinner!” Zaccheus reached his home, but he did not know what had happened to him or how he had climbed down the tree or how he had made his way through the throng. Breathless he entered his house; his eyes gleamed like stars in his flushed and dirty face. His wife hardly recognized him. She regarded him with mixed fear and wonder as he embraced her and stammered: “He—is—coming.—He—is—coming!” “Who?—Who?” “The Prophet—Jesus of Nazareth is coming to us!” “Then let me go. I must get things in readiness,” cried his wife as she broke from him. Should she not prepare the finest of food and decorate her home with the most costly of her possessions? But—no! Such display suddenly seemed shameful. Should she not rather hide all those things, gained by so much wrong? She hurried about the house, while Zaccheus looked about him in bewilderment. Then some one rapped, and Jesus entered. He stepped quickly to Zaccheus and took him by the hand. Forgotten were the beauty of the room, the torn clothes, the soiled hands. Everything—everything was forgotten; for beneath the gaze of those eyes he could feel but one thing—his tremendous debt,—the great penalty which he had loaded upon himself during all these years. And now Jesus had come to him—He did not despise him: No—he saw it all clearly—Jesus had forgiven him. He felt small and poor and rich and free at the same time—so free—so free! Jesus had forgiven him! Jesus had come to him! The burden of his life fell from him like heavy chains torn in two. No, he could no longer live as in the past. That was impossible. Then he began with stammering words: “Jesus—Master—the half of all my goods I give to the poor; and if I have taken any thing from any man by false accusation, I restore him fourfold.” Then Jesus pressed his hand and said: “This day is salvation come to this house.” “Salvation,” yes, truly. That which had tormented,—had seemed so impossible to him, was now all so easy, so beautiful. He was freed from his burden, because Jesus had come to him.—From E. and O. Zurhellen, \emph{Wie erzaehlen wir den Kindern die biblischen Geschichten}, pp. 322 ff.—N. B.! We are by no means of the opinion that Zurhellen’s example should be followed without modification; He takes so much for granted at which one may only guess; and he is not always simple and chaste enough; but, on the
\quellseite{580}
other hand, one can plainly see from the example what we (p. 474) said above about the motivation, the revealing of the acting person’s thoughts, and portrayal of the origin of the deed.

\subsection{3. PENETRATION. Zaccheus the Publican} (Recapitulation, comparison, and application are here combined with penetration.)

Of whom have we spoken in this story? Of Jesus and Zaccheus. We wish to pay the chief publican two imaginary visits. First we would visit Zaccheus before he knew Christ. We go to Jericho and ask for the house of Zaccheus. All know him, and any one can show the way. His house is larger and finer than the majority of houses in the town. The floors are covered with costly rugs. Along the walls are low, soft couches. Many servants come and go. All this costs money—much money. But money is plentiful with the chief publican. In a private room stands a large chest full of gold. Where did he get all this? We shall see that presently. Those under him, the publicans, come to him and he gives them his orders. “You must get money for me,” he commands, “much money. For every hundred-weight of grain that goes to Jerusalem I want a dollar, for every bushel of raisins, twenty cents, and for every yard of purple half a dollar.” He knew well that it was from two to four times as much as the emperor demanded; but he wanted to get ever more money, and therefore he did not care how much he demanded. And when the people complained and asked: “Is the tax really so much?”, his tax gatherers were instructed to answer, “Yes, so much.” What do you think of that? It was not right. That was lying and cheating. He took the people’s money in an unrighteous manner. Deception is as wrong as theft. The publicans, however, knew that they must obey; for if they brought their master too little money, they lost their job. So they oppressed the people, for in addition to what their master demanded, they took a little for themselves. How much scolding and cursing must have been done at the counter of the publicans. The people knew that they were being robbed. But what did Zaccheus care for their scolding if he got their money! He wanted to be rich.

Matters went favorably for the chief publican, and he became richer from day to day. But how was it with his heart? Certainly, not well. When he saw the poor little traders creep by with their thin, heavily laden asses, his conscience must surely have said: “Listen! It is not right that you oppress the poor in order to enrich yourself from the little that they have.” And when the pious Pharisees went by with looks of contempt, as if to say: “Miserable publican and sinner”,
\quellseite{581}
—that also hurt. So he had no true pleasure in his riches. Perhaps he could not even sleep in peace, because those whom he had robbed walked like ghosts before him. Thus you see that riches do not make one happy. What would be far better? A good conscience. But who could help him to get that? The Scribes and Pharisees despised and scorned him. If he had gone to their schools on the Sabbath to hear God’s word, he would have been thrown out. He would have been told: “You publican and sinner, cheat and rogue, enemy of your people and friend of the Romans,—get away from here! You have no business among good people, and our God has rejected you.” The Scribes and Pharisees were proud and hard and had no mercy. Poor rich Zaccheus! This is what we see on our first visit.

After a while we pay Zaccheus another visit. How different are things now! He is not the same man. His body has not grown; for he is the same little, fat fellow. But things have changed in his heart. One can see that in his face. The lines of care are gone from his brow and the sad look has disappeared. He looks calm and peaceful. He is friendly toward his servants. He sends for all who have ever been cheated by him or by his clerks. Ere long a brown camel driver comes to him. “Have my men ever taken too much money from you?” “They certainly have,—many a time.” “How much would the entire amount be?” “Hm, I think about fifty dollars.” Then the chief publican opens his desk and counts two hundred dollars before the wondering man. Why so much? That is four times as much as has been taken. “Is that satisfactory?” The man hardly knows what has happened to him. Thus one after the other comes and goes. The chest of gold has grown considerably lighter. Better look out, Zaccheus, else all your plans to become rich will fail. But Zaccheus acts as if he did not hear.

In the evening we behold a strange gathering before the house of Zaccheus. It is composed of unfortunates of all kinds, poor, lame, blind, cripples, etc. The chief publican who a while ago would scarcely have noticed them, now goes to each with kind words and gives them liberal gifts of money. In this way he distributes half of all his wealth. But will he not again force the money into his coffers? O no! He has strongly impressed upon his assistants that they are to exact not a penny more than the law demands. And with it all he is happier than he has ever been, and his happiness is written upon his face. What has happened to cause this wonderful change? is the question his fellow townsmen are unable to answer.

\quellseite{582}
We know the answer. The Lord Jesus alone has wrought the change. We have already told how kindly the Saviour treated him when He came to his house. He, the good, holy man in his, the sinner’s house; the love that shone forth from His eyes and never a word of censure or blame! It softened the publican’s heart. He knew that he had not deserved so much kindness and love. From all that the Saviour said and did, it became clear to him that He was holy, sinless. In the presence of Jesus Zaccheus saw himself for the first time as a great sinner, and all his deception and seeking after riches loomed before his eyes as evil. “What shall I do that I may become clean?” he asks. “I will do what I can; perhaps the Lord will forgive.” And then he said what we have already learned, “The half of my goods”, etc. You have seen that he was in earnest, and it happened as he desired. Certainly the Saviour forgave him: “This day is salvation come to this house.” Now the page has been turned. To be rich is no longer the greatest thing in his eyes; but his desire is to please the Saviour, to do what He wishes, and to hate what He hates. This the Saviour wrought by His love. Tell me the Scripture passage you have learned about His love. “Oh, taste and see that the Lord is good.” Zaccheus tasted that. “The Son of Man is come to seek and to save that which was lost.” Who else of his time experienced that? The publican Matthew.—And now we hardly recognize Zaccheus. He is a different man, a new creature. Old things have passed away. What, for example is past and gone? All things are become new. What do we see in him that is new?

\begin{center}
\begin{tabular}{ll}
\multicolumn{1}{c}{\emph{The Old:}} & \multicolumn{1}{c}{\emph{The New:}}\\
Seeking after riches, & Seeking to please the Saviour,\\
Love of gold, & Love to the Saviour,\\
Cheating, & Making good,\\
Hoarding, & Distributing to the poor,\\
Unfriendliness, & Friendliness,\\
Bad conscience & Forgiveness, a good conscience,\\
Sadness, discontent & Happiness and peace\\
\emph{is past.} & \emph{have taken its place.}
\end{tabular}
\end{center}

In Zaccheus has this word of Holy Writ been fulfilled: Therefore if any man be in Christ, he is a new creature (2 Cor. 5: 17). Let us learn this passage.

You can doubtless recall another man with whom all things became new. It was Saul, or Paul. What things became old with him? Hatred against Christ, persecution of the Christians. And what became
\quellseite{583}
new? Love to the Saviour and all believers, obedience to Him, consecration and work for Him. Who caused this change? Tell me what more you know about Saul’s conversion. Thus Paul learned the truth by his own experience—the truth which he so beautifully expressed in his letter to the Corinthians.—Even now Jesus makes new creatures of all who receive Him into their heart and home. There is much in our hearts also which does not please the Saviour. Our heart is by nature like a dark, dreary room. One does not know how ugly the room is until a light is brought in. So we do not understand the condition of our hearts until Jesus reveals it to us. Then we see all kinds of evil thoughts: Envy and hatred like Cain and Joseph’s brethren, and covetousness; for we constantly desire the best for ourselves. There are hateful creatures in the room of our hearts—spiders, insects, worms. The Lord shows us this by means of His word and the Holy Ghost. A pious man once said to two bad boys: “Pray for the Holy Spirit as often as the hour strikes.” They obeyed, and ere long saw what bad boys they had been. They continued to pray, and received hearts that loved the Saviour. We can not see the Saviour and can not, like Zaccheus, give Him our hand and say: “Welcome to our house, Lord Jesus!” But He is nevertheless here and sees and hears us if we mean it. But we must be in earnest. How did Zaccheus show that he was in earnest? He not only promised Jesus that he would be different, but he also kept his word. He did not do like many children who say: “I am sorry and will not do it again,” and within an hour forget their promise. They do not mean what they say.—If there is anyone here who has secretly taken something from someone else, or at home, he should do as Zaccheus did. He should at least confess, and if possible, give back that which he took, and not steal again. Otherwise he is not in earnest. How different things become in the home of a sinner, a drunkard, for instance, when Jesus enters there! Father, mother, and children are different in spirit and in outward appearance. There is even a change in the furniture, the walls and the windows. The heathen leopard becomes a lamb when Jesus works in the heart through the preaching of the Gospel. There is a wonderful change in a country when the people become Christians. A shipwrecked sailor was cast upon an island where cannibals used to live. In fear and trembling he climbed a hill and saw—a chapel. Then his fear fled, for he knew that where the cross is planted, he need not fear. Behold, I make all things new!—Fankhauser, page 141 ff. Naturally, the conversational form should here be made prominent.

\quellseite{584}
\subsection{4. APPLICATION. The anointing of David} Because God looks upon the heart, we can do nothing wiser than to strive for a clean, beautiful heart that pleases Him. It pleases Him best when we, like David, think much of Him and love Him. If you do not know how to do this, pray to the dear Lord. He will gladly give you that which He gave David—the Holy Spirit. He will teach you to love God and the Saviour aright. He will make you pleasing to God.—Certainly now and then evil thoughts will slink into your heart,—a hateful word, hateful thoughts against your companions, a foolish trick. But the Spirit of God tells your conscience kindly and gently that you did wrong. Would you remain pleasing to God? Then follow David’s example. How? Sing a hymn, repeat the passages of Scripture you have learned, or think of David the shepherd boy, or of Joseph. What did Joseph do when he was tempted to wrong? How did Joshua admonish his people? Therefore keep watch over your hearts, that you may please the Lord, your God.—You liked David’s courage, boys, did you not? I will tell you a story of a boy who had the right kind of courage. Story—The Proper Courage. Little Edward lacked courage to do wrong, to lie, to steal; but he was brave enough to save a child’s life at the risk of his own.—Above all, we want to equal David in one point more. Do not feel proud if you have been good and obedient, and do not boast how much you read your Bible or pray: only think, it is enough that the good Lord knows it. Many small and large children are like hens. If they have given a poor child some little thing, or if, for a change, they have obeyed their parents, or perhaps have not talked in school, they cackle loudly about it like a hen when she has laid an egg; or they look around to see if some one will not praise them. Do not do so. Rather think of the dear Lord. Rejoice if you can please Him. He sees it, and at the proper time it will be revealed.

The Lord no longer needs kings—yet He will make of us something finer than kings. He watches to see whom He can make His child and take into heaven, where each one will have far more glory than David ever possessed. But He will first conduct an examination. Each one must stand before Him to be tried, as the children of Jesse were tried by Samuel. But God will not only look upon us. He will look us through and through. On that great examination day all things will come to light that were hidden and secret, the good as well as the bad. Then perhaps a child will stand before Him and say: “Behold, dear Lord, here am I, Mary, the daughter of the
\quellseite{585}
rich miller. My father has much, much money, and all people say, You dear child, you pretty Mary! Dear Lord, do you not know me?” But God will sadly say: “I know you well, but I cannot use you. You have constantly thought of your pretty clothes and have forgotten Me entirely. You have kept yourself clean and nice, but you have not kept your heart and mind clear of evil thoughts. I look at the heart. Depart from Me”.—Or big John will say: “Dear Lord, I have always been the strongest in our school, and all have praised my strength and agility”. But God will perhaps say: “Yes, you have overcome others, but you have not conquered your pride and anger and the evil in your heart. Your heart is not good and you have never prayed to Me to cleanse it. Depart from Me!” Then will come clever Lydia, and say: “Dear Lord, do you not know that I have always been the first in my class? You will certainly give me first place in heaven, will you not?” But God will say: “Lydia, I not only take into account the Scripture passages and Bible stories they all can tell, but also this, that they practice what they have learned. But you have not obeyed your parents, and much less have you obeyed Me.”—Perhaps a small, shy boy will come and stand silently before the Lord. He was not strong, and the others often teased him. He was poor and had no shining shoes or warm cap, and was often cold in winter. But the good Lord will regard him kindly and say: “Come to Me, dear child. I have seen how you have thought of Me and loved Me. I know well how often you have shared your bread with other poor children and how you have overcome the temptation to lie or steal because you did not wish to make Me sad. I know you. You are My dear child. Enter into the eternal joy which I have prepared for you”—See how differently God looks at things! How easily could rich Mary, big John and clever Lydia have pleased God, if only they had kept watch over their hearts! Despise not the poor child in torn clothes. Perhaps God will be pleased with his heart rather than with yours. The Lord looks upon the heart. “Search me, O God, and know my heart: try me and know my thoughts: and see if there be any wicked way in me, and lead me in the way everlasting.”—Fankhauser, page 130 ff. The catechist would do better to confine himself to what Fankhauser says of the examination.

\subsection[5. A complete example]{5. A complete example: How God made the first people happy}

\subsubsection{Preparation} From whom does everything that we see on earth come?—Correct. It all comes from the dear Lord who made it in six days and still preserves it. Therefore we call God the Creator of all
\quellseite{586}
things and say on Sunday in church: “I believe in God the Father, Maker of heaven and earth.” But for whom has God made this world with all its beauty? What do you think, Anna, for whom has God made this earth with its mountains and valleys, rivers and seas, fields and meadows, birds and animals? Indeed, God made it all for us; we should live here; it should be our home and give us all that is necessary for our life. Therefore we should be glad and happy. How rich is our God, who made everything so beautiful; and how good and kind He is for calling all this into being for our sakes. And yet we have not heard the best that God made in the beginning, in order that we might be perfectly happy. What could that have been? We shall hear about it today, for we shall hear how God made the first people happy.

\subsubsection{Presentation} When God created the world, He made a most beautiful garden. You have perhaps seen pretty gardens, but none so beautiful as that. It was called Paradise, that is, Garden of Pleasure, because it was a pleasure to be in it and to live in it. There were beautiful trees of all kinds—forest trees which gave cool shade and fruit trees which yielded the finest of fruits. And Oh, the many, many flowers that grew there! They bloomed and gave out perfume as flowers never since have done. There were also all kinds of animals, large and small, from the little mouse to the lion, the king of beasts. Happy birds fluttered from twig to twig or flew high into the air and sang until the heavens resounded. There was plenty of water, for there were many springs, clear as crystal, flowing from the rocks. These flowed down the mountain-side and became brooks, and the brooks grew into rivers and streams. Four such streams flowed through Paradise and watered it, keeping the grass and trees and flowers green and fresh. Rich people, kings and emperors have made themselves fine gardens, but believe me, children, none of them, however large and beautiful, could compare with the pleasure garden which God made when He created the world. And think of it!—Adam was permitted to live in this beautiful garden of Paradise. Oh, how happy did God make our first parent!

He could go through the garden, rejoice over the flowers, eat the fruit and lie in the cool shadow of the trees. He could see all the animals and play with them. But with all this, something was yet lacking. What was it? He had no one to speak with. The animals followed him when he called them by the names he gave them and let him stroke them. The birds came and perched upon his shoulders
\quellseite{587}
or rested on his hand. But when Adam spoke to them they could not reply and they understood nothing of his thoughts. Then Adam wished for some one who would be his equal, with whom he could speak, and who could be a companion in his joy; for his heart was full to overflowing when he beheld all the beauty that surrounded him. The good Lord, understanding Adam’s wishes, let him fall into a deep sleep. Then He took one of Adam’s ribs and out of it made a woman, Eve. When Adam awoke, he saw her stand before him. He beheld her with wonder and shouted with joy: “This is now bone of my bones and flesh of my flesh: she is a fitting companion and just what I have wished for. I can speak to her and she to me. I can tell her all my thoughts.” How happy was Adam then! He went about the garden with Eve and showed her a pretty spot here and another there. They sat beside the brook and listened to the babbling water, or enjoyed the sweet songs of the birds. And when they became hungry, they plucked berries or apples and pears from the trees. But many a time they must have stood still and said to each other: “Oh, how much our God must love us since He has made everything so beautiful and placed us in this lovely garden.” They worked, too, for God had told them to take care of the garden. They watered the tiny plants, making little ditches, so that the water could flow to them; they lifted up a branch of the grape-vine here one and there one and fastened them so that they could climb up the tree. But their work was not hard, it did not make them tired. It was a pleasure and a joy and only increased their happiness.

God had made Adam and Eve masters not only over the garden, but also over all animals that were in it, the small as well as the great. “Be fruitful and multiply,” He had said, “and replenish the earth, and subdue it; and have dominion over the fish of the sea, and over the fowl of the air, and over every living thing that moveth upon the earth”. In our picture the artist has shown us the Lord God pointing to the animals with one hand, while He raises the other before Adam and Eve and speaks these words to them. See, here is a mighty elephant with his big ears and long trunk. Above him sit two mighty birds. Farther down is a proud horse with high head and flowing mane. And near by—what is that just showing its little head?—Sure, it is a lamb. And that other is a deer, which can run so swiftly and flees when it sees human beings. What comes next? Correct—that is a lion. The lion is a powerful animal, and therefore he is called the king of beasts. Beside him a cow is peacefully lying,
\quellseite{588}
while between her legs a serpent crawls. Could we today put all the animals together as the painter has done? Why not? Certainly. They would kill and eat each other and the people, too. We would have to put some of them behind iron bars. But then it was different, for when God created them the animals were not enemies to each other. It was easy, too, for Adam and Eve to be their masters, and it was a part of their happiness when God said to them: “You shall rule over all these creatures and over the whole earth and subdue them”. So God made the first people kings over the earth.

But, my dear children, I have not yet mentioned the best and greatest happiness of Adam and Eve. What could that be? It was when the Lord God came to them in the garden and spoke with them as a father with his children. Oh, how they rejoiced when He entered the garden, and how they hastened to meet Him, as children do when their parents return from a journey. How bright their eyes!—They had so much to show Him and so many questions to ask. Today we can not see God, but then He came to them in the form of a human being, so that they could see Him and feel Him. What delightful conversations those must have been between God and them! He would explain to them what they did not understand, how He had made the world, that it was all for them for their enjoyment and happiness and that they should rule over it. “Yes,” said God to them one day, “everything you see is yours and you may eat of every tree in the garden, but of this tree here, the tree of knowledge of good and evil, you shall not eat. For the day you eat thereof, you shall surely die.” What a kind God, who, with this one exception, placed everything into their hands. How He must have loved these first people, and how happy they must have been!

\subsubsection{Penetration} (association, comparison, and abstraction). Where did the first people live? Where was Paradise? What—you do not know? Neither do I. Many learned men have thought that it was in Asia Minor, between the Euphrates and Tigris rivers. (Show map.) They may be right, but no one knows for sure, and we can not seek it, for it is no longer upon earth. Since when has it disappeared from the earth?—At any rate, since our first parents fell into sin there has been no Paradise here below. Much more important is another question, namely, how Paradise looked and what it was like. And to this the name “Paradise” gives the answer. What is the meaning of the word Paradise? Who has
\quellseite{589}
noted it? Paradise, or Garden of Eden, means garden of pleasure. There are pleasure gardens today, for many a rich man has one made for himself and many cities have parks which could be called pleasure gardens. But why call them pleasure gardens? Because the people find pleasure and enjoyment in them. When they become weary, they can rest there. When the sun is hot, what can they find there? Especially who can play and have a good time there? You children, of course. But all the pretty parks and gardens, with which one can they not compare? Now tell me all that was to be found in the pleasure garden called Paradise. What was there besides the beautiful flowers and the trees? What can you tell about the streams that flowed through and watered the garden? Because of this, what was the condition of the grass and all plants and trees in this garden? It is often so different now. What happens so often in summer to the grass and plants? Yes, sometimes the trees lose their leaves in summer when it has been hot and dry and no rain or dew falls. It was not so in Paradise. Everything was as fresh and green as in spring-time. And just think of it, children, there was no hot summer to make everything dry; no autumn when the leaves fall from the trees; and no winter, when all nature seems dead and snow and ice cover all the earth. To be more correct, these four seasons existed also in Paradise, but none of them brought those disadvantages they bring now. Oh, then, what pleasure and joy to live in that garden! Therefore, how do we rightly call it? And whom did God place in that garden to live? What should Adam according to God’s will have there? Surely, pleasure, satisfaction and joy. Truly happy did God want to make Adam, therefore He placed him into this garden. Let us all say this together: \emph{God wanted to make the first human being happy, therefore He placed him into Paradise}. (Write on blackboard.)

Was Adam happy in Paradise? Tell us about his happiness. What, besides the trees and the flowers gave him the most pleasure? To be sure, the birds and the other animals; they had life and could look at him. Tell me how trustful they were toward him. What did he give to each bird and animal? Adam had closely noted their way of living, everything was so new to him, and God had given him keen understanding. If we now knew the names which Adam gave to the animals we would certainly wonder at their fitness. The animals in time learned the names Adam gave them.
\quellseite{590}
How is this shown? But when the dog came at Adam’s call, and the lion lay at his feet, so that Adam could fondle his mane, and when he spoke to them and called them pet names, what could they not do? How do we feel without some one to whom we can speak? Yes, we feel lonesome and forsaken. Just fancy yourself alone, with no one to speak to. How unhappy you would be. How would Adam therefore have felt in spite of the beautiful garden in which he lived? He would certainly have been without something very necessary to his happiness. Man is happy only when he can express himself and finds that his thoughts are understood and answered. The reason why it is so hard in prison is not, because one must work, but because one must be alone and dare not speak to others. What, then, did Adam soon wish for? To whom was his wish satisfactory? What did God say to Himself? What did God immediately do with Adam’s wish? Why would He act upon it immediately? Certainly, He wanted Adam to be happy. Could He not have given Eve to him at the very beginning? He could have done so, but why did He not do it? If no one knows, I shall have to tell you. The Lord wanted Adam to learn what he yet needed. If we have something from the beginning, then we do not yearn for it or value it highly. On the other hand, if we receive something which we need and for which we have long wished, we are much happier when we get it. So it should be with Adam. Tell me how God gave him some one to whom he could speak. How did Adam feel when he awoke from his sleep and suddenly saw Eve before him? How did he express his wonder? He said with these words that here at last was a suitable companion. Where, so far, had he found no suitable companion? It was immediately clear to him that man and woman belong together. With what wonderful words did he show this? What do we say of a man and woman when they leave father and mother and live together? And by what name do we call that life when a man and woman are united? You do not know it, and therefore I will tell you: They are married, and such living together we call married life. Who gave Adam his wife? God Himself thereby instituted the marriage relation and showed that it is pleasing to Him when man and woman enter into the same. If God Himself instituted marriage, what should husband or wife not do?—What is done with the marriage relation when husband or wife separate or do not love each other and are not true to each other?—Yes, they break the marriage tie,
\quellseite{591}
for they break the word, or promise, they have given to each other. That is a great sin. Tell me the Sixth Commandment.—But we must return to our story! Did God make Adam happy when He gave Eve to him? And now Adam had some one with whom he could do what? And what could Eve do when Adam spoke to her? Tell me in a few words how happy they were. Adam and Eve could then say what King David later declared. When he thought of all that God had given him to make him happy, he rejoiced in these words: “The Lord is my shepherd; I shall not want”, etc. Did Adam and Eve have to work? But why did this work not disturb their happiness? Correct. Joyful work is a pleasure and no burden. Now, who gave Eve to Adam, thereby making him happy? Let us repeat together: \emph{God made Adam happy, because He gave Eve to him, with whom he could speak and rejoice}. (Write this on the blackboard.) Mary, you may tell us this. Now Frank may repeat it. John, what was the first thing we heard that God did to make Adam happy? Yes, God made Adam happy when He placed him in the pleasure garden. Louise, what did we hear further?

However, we know still more about the happiness of the first people. God not only placed Adam in Paradise, gave him Eve to be his playmate and companion, but also made them rulers over the animals in the garden and over the whole earth. Who can repeat the words with which God gave them this right? Yes, those are the words. Many people should come from Adam and Eve, so that the entire earth should be filled. And wherever these people should go, they should rule over all the animals, and all the powers of the earth should serve them. Tell me, Louis, how do you think the people should make the beasts serve them? How do we today compel the horse to serve us? How does the cow serve us? The chickens? So should all animals, even the lion and the elephant, serve Adam and Eve, each one according to his powers. I have also spoken of the powers of nature. What had I in mind? Is there none who knows? I was thinking of the air, the wind, gas, electricity. Are not those great powers with which man accomplishes wonders? How do we today see that we control the air? How do we use the wind? For what purposes do we use gas and electricity? And now think of the treasures beneath the earth. Can you mention some of them? Certainly, coal, iron, and the other metals, silver and gold. The good Lord set Adam and Eve
\quellseite{592}
over all these. They should use them all for their comfort, good, and pleasure. How rich God made them! The lords and kings of the earth also have control over many animals, powers and hidden treasures, but they rule over only one land, while Adam and Eve were rulers over the entire earth. What title then can we rightly give them? Perfectly correct. We can call them kings, for they were kings and rulers over the whole earth, with all that was in or upon it. People often follow an earthly king because they must and because they will be punished if they do not. How did the animals and the powers of nature obey Adam? And what made it easier for him to control the animals? Recall the picture you saw the last time. Which beasts did we see peacefully lying side by side? Could that be done today? Why not? What must be done with the lion or the serpent? They would not only hurt the other animals, but whom also would they attack? Yes, man's life would not be secure. Then what was not present among the animals of that time? And between whom also was there no enmity? That is right. As the flowers, large and small, did not wither or fade, so was there no enmity, no death among the animals. They did not prey on each other, but all lived together in peace and harmony. It was good to be king and ruler there. Who made the first people rulers over the whole earth? How did He want to make them? Yes, He desired to make them happy. Let us say together: \emph{God made the first people happy, because He made them rulers over the whole earth}. (Write on the blackboard.) In how many ways have we thus far learned that God made the first people happy? Tell me the three ways.

Was there anything more that God could do to make the first people happy? One might think that they had all that was necessary to make them happy, and yet the most important thing was lacking. Picture yourselves living in the most beautiful garden there ever was, where you had brothers and sisters to speak and play with; where you were rich and ruled over all things, but where you had no father or mother. Would you be really and truly happy? And how would it be if father and mother were living, but they were angry with you and wanted to have nothing to do with you? How would you then feel? Certainly you would be very unhappy. Of whom would you be constantly thinking? And of what especially would you be thinking? If you were good children, all joy would depart if you had to say: Father and mother
\quellseite{593}
want to have nothing to do with us. They are angry with us. Now, who was Adam’s father? Why was God Adam’s and Eve’s father? How would the first people certainly have been, if their heavenly father, the dear Lord, had never come to them and wanted to have nothing to do with them? What, then, would God have done to make the first people perfectly happy? Correct. God came to them in the garden and spoke to them as a father with his children. Could they see God? Yes, children, I believe they could see Him as we some time shall behold God in heaven. What did God have to do, if Adam and Eve were to see Him with their natural eyes? Of course, God must make Himself visible. God is a spirit, but He can take on human form. And what happiness must have been theirs when Adam and Eve saw God come to them in human form. What do you think they did when they saw Him coming to them? And what do you think they did when He walked through the garden with them? Tell me a few questions which they probably asked Him? There were many wonderful things in this beautiful garden which they could not explain. Which was the first question, in your opinion, they asked? Yes, I too, think that they asked the Lord who made the beautiful garden, the whole earth, and the heavens above. O children, those must have been blessed hours when God spake with them and they told Him all that filled their hearts. Now what could not but be impressed upon them as they talked with God? It was out of love to them that God made everything so beautiful; out of love He came to them and communed with them like a father with his children. And I should greatly wonder if Adam and Eve, when they were alone, did not often say to each other: How dearly God must love us! Oh, we can not do otherwise than love Him in return. Of what Scripture passage does that remind you? Yes, so would they have said to each other: “We love Him, because He first loved us”. Now we know what made Adam and Eve happiest. What was it? Let us repeat together: \emph{God made the first people happy, because He came to them and walked with them as a father with his children}. (Write this on the blackboard.) What was the first thing that God did to make Adam and Eve happy? What was the second, the third, the fourth? Yes, God had done so much for them, and it is no wonder that they were happy, perfectly happy. And why did He do so? Right. He did all this purely out of love to them.

\quellseite{594}
And notice, children, this also, that it was because of His love that He gave them only one commandment. He did not wish to plague them with many commandments. They should not constantly have to say: Yes, God has forbidden this and that. What is the commandment God gave them? \emph{You may tell me—and you—and you}. Why did God call this tree the tree of the knowledge of good and evil? Do you not know? Well, the word “knowledge” comes from the word “know”. This tree should help them to know what was good and what was evil. Does it make any difference whether we know what is good and what is evil? Yes, indeed, much depends upon it. You only need to compare yourself with father and mother to see it. They know much better than you what is good and what is bad. They have often enough experienced how bad it is when one does wrong, and how blessed it is when one does right. Hence they do not easily permit themselves to be persuaded to do wrong. And was it not necessary for Adam and Eve also to learn what was good and what was evil, so that they might always choose the good? And what was to be the result if they did the good? Correct. Their happiness was to increase. Are you happy and joyful when you have done wrong? How do you feel? Yes, you feel ashamed and uneasy. It should not be so with Adam and Eve. \emph{They should learn most precisely what was good and what was evil, in order that they might always do good and be happy and blessed}. (Write on the blackboard.) That was what God wanted, and therefore He placed the tree with the forbidden fruit in the garden. What, therefore, were God’s intentions in regard to them? His intentions were good because He loved them. Now one should have expected that Adam and Eve would have done what? Yes, they should have loved God. They said many times that they loved God; however, they should not only say it, but also prove it by their deeds. How could they have best proved that they loved God? Perfectly correct. That is the best proof of our love to God, that we keep His commandments. In which Scripture passage did you learn that? We shall see if Adam and Eve loved God so well that they kept His commandments. That would surely have been the best thanksgiving to God for having made them so happy.

\subsubsection{Application} Children, does God love you and has He made you happy? At any rate, where has He not placed you, that you might live there? Why would that not do? But is not this world in which
\quellseite{595}
you are permitted to live a beautiful one? When do you notice most that this world is beautiful and glorious? In spring-time, of course, when everything blossoms and grows. It is sometimes so lovely that we hear the people say: “It is like heaven.” And as Paradise once was the most beautiful spot on the earth, so, I think, God has also assigned to you a place upon this wide earth where you have it best. Which place do I mean? I am thinking of your home. That is for good children the best place. Ask the grown-ups, who have been among strangers for a while, or those who have no home, if there is a better place on earth than home. It is true, one can be richer and freer elsewhere; but, as a rule, it is nowhere better than where father and mother are. Think of the prodigal son. What did he think at first? What did he learn afterwards? And after what did he finally yearn with his whole heart? Yes, father’s hand holds closest, mother’s eyes beam kindliest, and mother’s heart understands best. Now, who placed you in this beautiful world and gave you father and mother and home? What state of mind should this work in you? Did He leave you without any one to speak to? You certainly would be unhappy children if you had no one with whom you could talk, for it is very difficult for you to sit quiet and silent for but an hour. But whom did God give you to speak to at your pleasure? And who is it that always listens to you and never becomes weary of your prattle? See there, how happy God has made you! and when you become older, you will more and more understand that something of that kingly position which God once bestowed upon Adam and Eve, still remains for you. And if you had no home, no father, no mother, no brothers and sisters or friends, why would you still not be forsaken? Yes, indeed, God in heaven would still be your father and you would be His children; and that is the greatest happiness that one can have. Through what has God become your father and you His children? In which passage of Holy Writ have you learned this? Correct. The great thing in Baptism is that it makes us God’s children. What does a true father do for his child? What does God therefore do for you, because He is your father and you are His children? How does He care for you? In what other ways? And still others? We would not forget the most important thing. The dear Lord sees to it that you hear His word and that therein you are shown the way to true happiness. He therefore provides not only for your body, but also for what? He would make you
\quellseite{596}
happy in body and soul. I know a hymn in which it is beautifully told how good we have it with the good God, with our dear Saviour Jesus Christ. All of you know it; we have sung it often together. It speaks of the Lord Jesus as the shepherd and of us as His lambs. Now you know it; repeat it, John. But the finest thing is that you as baptized children can also speak with God like children speak with their father. How may you call God? Who has taught us that we can call God by the name—father? Where? What are the opening words of the prayer Jesus taught us? How does our Catechism explain these words? God would give us encouragement and pleasure in praying. That is why He tells us especially to call Him father, for His words are a tender invitation to pray to Him. How should we pray to Him? How do children ask their father? Yes, they are not timid and fearful, but what confidence do they place in him? Now is it not something unspeakably great to be able to pray to God in heaven with all confidence and assurance as dear children ask of their dear father? We must call that true happiness. Who has given you the right to speak to God at all times and tell Him all you think and need? How has God then also made you? Let us say it together. \emph{God has made also us happy, for He has made us His children through Baptism}, who can entreat Him as dear children entreat their dear father. And, children, whatever may be lacking to make our happiness complete, He will give us above in heaven; for if we die His children, He will take us home into the heavenly paradise. There our happiness will be full. The last stanza of a certain hymn learned by us recently, fits in here. Will you say it? How must God be disposed toward you if He showers so many blessings upon you? How should you feel in your hearts since God has made you so happy? What kind of children would you be, if you did not feel so? Yes, indeed, you would be ungrateful to God, you would be thankless, evil children. And how must you show that you love God? Yes, that is love to God when we keep His commandments; and the better we keep them the greater will our happiness be.

\subsection[6. Presentation of the Bible History]{6. Presentation of the Bible History in the Sunday School}

(Here the most important parts of the penetration and application must be combined with presentation):

\subsubsection{How the first people sinned}

\emph{Preparation}. How happy God made the first people! He let Adam live in Paradise, and Paradise was the most beautiful garden the earth has ever seen. Here Adam could have all that his heart
\quellseite{597}
desired. Here God created Eve for him, that he might have some one with whom he could speak and share his joy. Oh, how happy both of them were when they wandered through the garden hand in hand and called each other’s attention to the beautiful nooks they had found. I wish that I, too, might have been there. Then, besides, they were masters over all that was in the garden—over the beautiful flowers, which they might pluck as they wished; over all the sweet fruits, which they might eat as they desired; over the birds of the air, which fluttered about them and showed no sign of fear; over all the beasts of the field, even the mighty lions and giant elephants, which obeyed their slightest word. Yes, God made Adam and Eve rulers over the entire earth. Indeed, they were happy people. And yet, you know that their greatest happiness consisted in something higher, much higher. It consisted in this that God came to them in the garden and spoke with them as a father with his children. What blessed hours those must have been when they so freely and joyfully associated with God as you children do with your parents, and when without fear of the almighty and holy One they could look into His eyes and walk at His side! He did not burden them with many commands, so that they had to watch every step, lest they transgress one of His commandments. He gave them a single command, and that was for their good. What was that commandment, Frank? That is correct. They might eat of all the trees in the garden save one. It should be both church and school for them, for from it they should learn what was good and evil, that they might always do the good and thereby grow in happiness and blessedness. For what was God’s purpose in all that He did for them? That they might be truly happy.

What should these first people have done with this command of God, Mary? Yes, they should have gladly kept it. And if they had kept it, their happiness would never have ended—it would have lasted for ever. But they did not keep it. Today we shall learn how the first people disobeyed God’s command, how they sinned.

\emph{Presentation}. How may it have happened that the first people transgressed God’s command and sinned against Him when He had nothing in mind but their happiness? One would think that it could not have been possible for them to sin against God, so well had He arranged all for them. At first, Adam and Eve did keep God’s command. How long that lasted, we do not know; but it
\quellseite{598}
could not have been such a very short time. More than once they thought of the command and said to one another: Of that tree yonder we may not eat, lest we die. I think that when they looked at the forbidden tree, a sort of horror stole over them. I should not wonder but that they came after a while to call the tree “the tree of death” and passed it at a distance rather than go near and touch its fruit. They evidently thought not only of God’s command and the death penalty that would follow disobedience, but also of His love. And they probably said to one another: Since God loves us so, and has made us so happy, we will love Him in return and try our best to keep His command. But how did it come that they finally transgressed and ate of the forbidden tree?

The evil thought did not spring from their hearts. But there was one who could not bear to see them love God and obey His command. That was the devil. He is the great enemy and adversary of God who is always striving to destroy God’s works. As often as he saw these people, so pure and holy and obedient to God’s will, he became angry and planned how to destroy this beautiful work of God, rob the pair of their happiness, make them impure and unholy,—yes, as evil as he himself was. There are boys who, when they are dirty, like nothing better than to make others unclean also. It is just so with the devil. He wanted to make these first people as evil and unclean as himself, so that God could no longer use them, could no longer be their Father, but would have to cast them out of Paradise. He said to himself: If I have once torn them from God, then they are mine and will serve me all their lives; then I shall always be their master and they shall ever be my slaves.

The devil evidently considered long how he might accomplish his purpose. He asked himself: How shall I begin? How can I lead Adam and Eve astray, so that they disobey God and become mine forever? It was clear to him that he could gain his end more easily and surely if he tried Eve instead of Adam; for the woman, being weaker than the man, could more easily be led astray. He also told himself: As I am, I may not show myself. I dare in no way betray that I am the enemy of God and their enemy. I must appear to be their friend who is seeking their best interests. Finally his plan was perfected. The devil knew what he would do.

Craftily and cunningly he began. How craftily and cunningly, the Bible tells us at length. In Paradise there were all kinds of
\quellseite{599}
serpents, large and small. They glided swiftly upon the ground; they coiled themselves in the bushes or wound themselves about the trunks of the trees and so made their way to the branches. Adam and Eve were not afraid of them; for they could not and dared not harm them. They had frequently watched them as they played with the other animals, and wondered at their cunning, in which they surpassed the others by far. It was one of these serpents that the devil used. It should be his mouthpiece, the agent through whom he would deceive Eve and bring her to transgress God’s command.

Eve generally avoided the forbidden tree, but one day, as she drew near to it, she saw an especially pretty serpent coiled about the trunk of the tree and looking in her direction. It seemed to Eve that the serpent was looking directly and steadily at her. She stood there and gazed at the serpent and the tree. And as she looked, the serpent suddenly began to speak. Eve was filled with wonderment, for though she had often before seen the serpent, she had never heard it speak. She thought that no one could speak like she, her husband and the good Lord. Can serpents speak, Mary? That is right; they can not speak. Then how did it happen that the serpent on the forbidden tree could speak, Louise? We know the answer: it was really not the serpent, but the devil who spoke through it. Eve should not know who was speaking to her, therefore the devil spoke through the serpent. And what did he want to do? Who remembers? You are right, John. The devil wanted to deceive Eve and induce her to transgress God’s command.

And now hear what the devil said to Eve: “Did God really say, You shall not eat of every tree of the garden?” He knew well that God had said something different, namely, that God had given Adam and Eve permission to eat of all the trees with but one exception; but he wanted to find out if Eve still remembered God’s command. He wanted to make her uncertain about it. He said to himself: If she is no longer sure of what God told her, I can the more easily get her to disobey and sin against God. We already begin to see how cunningly he started to carry out his plans. But this time his trickery did not help him; for Eve knew what God had commanded, and was not in the least deceived. She replied: “We may eat of the fruit of the trees of the garden; but of the fruit of the tree which is in the midst of the garden, God hath said, Ye shall not eat of it, neither shall ye touch it, lest ye die.” She evidently meant, What kind of foolishness are you speaking; the very opposite is true. God has by no means
\quellseite{600}
said that we may eat of none of the trees of the garden. He is not so cruel that He would put us in this garden with all the many fine fruit trees, and then tell us that we should not eat of them. He is too loving and good for that. He plainly told us that we might eat of all the trees of the garden except this tree about which you have coiled yourself. The dear Lord has His good reasons for making an exception of this tree. Whoever touches the tree or eats of its fruit must die. God does not want us to die. He wants us to go on living and to remain in this beautiful garden forever. What is the matter with you, you foolish serpent? You are generally so clever and know more than all the other animals.

In this manner Eve answered the serpent and the devil who spoke through it. The devil might then have left, for he had lost his case. In spite of all his craftiness, Eve had put him to shame. She not only knew the command exactly, but she also knew that God had given it for her good. But the devil did not consider himself beaten. If he cannot accomplish his purpose one way, he tries another. So he thought quickly and formed another plan by which he could gain his end and bring Eve to disobey. Stop, he said to himself, Eve thinks that God gave the command out of love, because He is seeking her highest good. This faith I must take from her. I must induce her to think that God gave this command out of envy and not out of love—not because He means well with her, but because He does not wish her to have the best. Therefore he quickly replied to the woman: “You shall surely not die if you eat of the fruit of this tree. God has lied. He well knows that the day you eat of it, you shall become like God, knowing good and evil. God does not want you to be like Him, therefore He gave that command. He did so not out of love, but out of envy and jealousy.” What the devil told her was a big lie, and he knew that he lied. But he did not care for that, if only he could bring the woman to believe him and eat of the forbidden fruit.

Did she see the lie? She should have done so; for she knew God; she had been with Him long enough, and had often enough beheld how holy and truthful He was and how He loved her with all His heart. She should have replied to the devil: Depart from me. I will not listen to you. God is not as you say. Would He have placed Adam and me in this beautiful garden and have made us so happy if He had been envious and jealous of us and instead of wishing us well? Thus she should have answered and thereby have smitten the devil a second time and have remained God's dear child. Such thoughts evidently arose in her mind, but other thoughts were now becoming
\quellseite{601}
active. The serpent had told her something she could not forget. It was this: You shall be as God. She thought that must be something great, to be like God and to know both good and evil. And the more she thought of this the more her present happiness did not seem worthwhile. It was with her as sometimes occurs with you children. Bread and butter tastes very good to you until you see a nice, fresh cake on the table; then you suddenly become dissatisfied with the bread and you want the cake. Or, how satisfied you are to be in the house with mother, until you see the children play outside! Then the room suddenly becomes too narrow and small, and you think: We should be really happy if we could play outside. It was something like that in the case of Eve. Paradise was suddenly not good enough. She wanted more—to be like God, and like Him to know good and evil. She forgot entirely to ask if Satan was telling the truth or if he was lying again. She forgot that God is certainly not envious and that He cannot lie. She heard and thought only one thing: to be like God. Her whole heart was suddenly full of that desire. But that was an evil desire; it came into her heart from the devil and not from God. But she did not ask about good and evil, she yearned after but one thing—to be like God. And then, children, Eve lifted up her eyes and looked at the fruit of the forbidden tree. Oh, if she had only not looked, she might still have considered the right and not have given her will over to evil desire. But she looked up and examined the fruit on all sides, and because it looked so good and sweet and promised to taste so good; because she believed the devil’s word that it would make her like God, think, O children! she finally stretched forth her hand, took the forbidden fruit, ate of it, and gave of it to her husband. O children, the angels in heaven must certainly have wept, and it must have cut God to the heart; for now the first people had believed the devil more than God! Now the command was transgressed and the first sin committed upon earth! Now were come upon earth death and all manner of misfortune. “By one man sin entered into the world, and death by sin; and so death passed upon all men, for that all have sinned.” The devil was the cause, for he deceived the first people; but they were not blameless, because they believed the devil more than God. If they had resisted the devil, he would have fled, but they listened to him and did his will. Therefore: Be not deceived! Evil communications corrupt good manners. Resist the devil and he will flee from you. Children, when evil companions entice you, follow them not. There are many sad days, but the saddest of all, as dark as midnight, is that day on which the first people sinned.

\quellseite{602}
\subsection[7. Review and drill]{7. Review and drill at the end of a longer section:}

\emph{At the conclusion of the history of Abraham’s life. A.) Bible story: 1. Headings of main and subordinate divisions of history.—2. Character sketch of Abraham:} a) toward God: pious, god-fearing, obedient, believing; he fears, loves, trusts God above all things; is grateful and humble. b) toward men: peaceable, unselfish, merciful, helpful, courageous, valiant, noble, hospitable, polite, friendly; an interceder, a faithful husband, a providing father, a good, kind master to his servants, a faithful head of the home who teaches God’s word to all his household. \emph{3. The promises given him}, which were already fulfilled; the three great promises, great nation, land, blessing to all people.—\emph{4. Character sketch of Eliezer.—5. Character sketch of Sarah and Rebekah.—B) Doctrine}: \emph{1. God’s essence and attributes}: almighty, truthful, faithful, long-suffering, patient, merciful, but also righteous, jealous, holy, impatient of evil. \emph{2. Human Virtues and Weaknesses}.—Which virtues did Abraham exhibit toward God and men? Which Eliezer, Rebekah? Which did Lot show; which Sarah? On what occasion did they show them? What should a master, the head of a family, a husband, a wife, a servant, be like?—Especially the conceptions of idolatry, unselfishness, magnanimity, worldliness, and immodesty, still largely alien to the child mind, should be diligently reviewed and explained. \emph{3. Religious concepts}: Altar, covenant, oath. \emph{4. Parts of Catechism as derived from the several Bible stories. 5. Stanzas of hymns. 6. A review of the geographical material and of that pertaining to cultural history should not be forgotten}.—From F. Zange, \emph{Leitfaden fuer den ev. Religionsunterricht}, vol. i and ii, 1906. A simpler arrangement would be to assign central position to the subject: Abraham the Father of the Faithful, and to group all material around it.

As a basis for a review of the entire history of Israel, Isaiah 5: 1–7, is particularly adapted; the class should show in detail 1. What God did for His people since their deliverance out of Egypt; 2. How sour instead of sweet grapes grew on the vine of His planting; 3. How the Israelites incurred His wrath.

\subsection[8. Insertion of the geographical element]{8. Insertion of the geographical element and of that of cultural history at the proper place for a broader survey}

a) Israel entered into Canaan and the Lord gave His people rest from all their enemies and showed them all the kindness He had promised. The teacher adds the following: The Israelites fared very well in their new land. The shepherds in the land east of the Jordan and in the
\quellseite{603}
mountains of Judah lived in tents, but among the rest of the Israelites each family had its own house, and an arbor of grape vines around it. The vines grew rank and luxuriantly, climbing all over the arbor and covering the flat roof of the house. In summer and autumn there was a roof of grapes from the arbor to the house. What advantage was there in that for the people? (Plenty of shade.) It was their delight to spend their time there in the shade. Only when they could not remain outside, did they go into the house. When was that? (At night and when it rained.) Of course, their houses were not so comfortably arranged as ours; for they had only one room, with one door and no windows. You already know what kind of roof they had. What kind was it? You would not care to live in such houses, but the Israelites were well satisfied. Why? They knew of no other kind, and they needed nothing different, for they could be outside all day.—Summary, House.

It was delightful out of doors. About the towns and villages stretched the \emph{fields and meadows}. What grew on the meadows? What grew in the fields? (Wheat and barley.) On the hillsides were forests of oak and pine trees, groves of palms, orange, and citron gardens. What fruit was there in profusion? Besides, there were fine apples, plums, apricots, peaches, almonds, and the like. There also grew the fruit of which the spies brought samples. Do you recall what they were? Already in January, when we have snow and ice, the almond tree was blooming there, and the fields were covered with violets, tulips, lilies, etc. At the end of April the barley harvest began, and at the end of May, the wheat harvest. Then came the fig harvest, and from August on the grape and olive harvest. You know what olive oil is; for what purpose is it used? This oil is made from olives. Olives grow on trees about the size and color of our plum trees, and are of different shades. There are blue, black, red, white and green olives. They are pressed, and the oil that comes from them is valuable. Olive-trees are also called oil-trees. Upon what mountain to which Jesus often went with His disciples were such trees? From what did this mountain get its name?—Summary, fruits of that land.

The Israelites \emph{worked diligently} in this beautiful land. Seedtime came in October. What did they do then? That was done as soon as the first rain fell. From the end of October until the beginning of March there was much rain. What time of the year is that with us? But there was also \emph{plenty} to do during the rainy
\quellseite{604}
season, especially when heavy rains damaged the vine-yards. What damage? (Ground washed away, vines washed out, vine-yard walls undermined and thrown down, etc.) What did they then have to do?—The last part of March and the beginning of April was the short rainy period. This they called the late rain. After that the sun shone hot from the cloudless sky, and the grain soon ripened. Which ripened first? Then began a season of continual labor. What had to be done with the ripe grain? Where do we put the sheaves from the field? The Israelites had at that time no barns, and so they threshed the grain upon large flat places in the fields, which they called threshing floors. These floors were circular in shape and were surrounded by a low wall. How do we thresh the grain? With flails and threshing-machines. The Israelites filled the threshing floor with grain and then drove oxen and cows around over it all day. How did the grain get out of the ears? (It was trodden out by the animals.) Then the grain had to be separated from the straw. In the evening, when the sea wind blew, the men took a shovel and threw the trodden grain in the air. What then happened with the heavy grain? (It fell upon the floor.) And what happened to the lighter straw? (It was blown away.)—Summary, Seed-time and harvest.

\emph{Harvesting} was joyful work. By what do we know this? (Their singing.) At this season all the people remained in the fields day and night, masters as well as laborers. Why? (To guard against thieves and fire.) They also prepared their meals there; but, naturally, they were very simple. Ears of grain were roasted at the fire, then rubbed out and eaten. Or the grain was ground between stones. What was made of the grain? (Flour, really groats.) The meal was then mixed with water into dough and thin cakes were then baked of it. And how did they quench their thirst? (With water.) They mixed the water with vinegar. Why? (Such a drink quenches thirst much better.) Where did they carry the grain? (Into barns.) There was little wood in the land of Canaan, and to build barns was too costly. So they made great caves in the earth.—Summary, Harvest.

The \emph{grape harvest} was still more joyful. During the grape season, from August to November, the entire family lived in the vineyard. In nearly every vineyard there was a tower. Why? (To stay at night and in case of rain.) Why did the people stay so long in the vineyard? (There was much work to be done there.)

\quellseite{605}
The grapes grew much more luxuriantly there than they do for us. Many a grape-vine was as thick as a tree and bore grapes without number. Such a vine bore as many as a thousand clusters of grapes. What did they make out of so many grapes? (They cooked them into grape honey and prepared them for wine.) The wine they made in the wine-press. That was a place sunk in the ground, or hewn out of a rock, about 2 yards long, 2 yards wide and a half yard deep. (Show the size.) The men tramped upon the grapes in this wine-press. What flowed out of them? (The juice.) Before every wine-press there was a deep place hewn out of the rock, and a trench led from the press to this basin. In the press and at the mouth of the trench there was a grating. Why? (To serve as a sieve and hold back the grape-skins.) What were the Israelites as long as they were in the vineyards? (Vine-dressers or wine-growers.) Why were these wine-growers so happy at their work? (Good harvest, juicy grapes, sweet must.) How did they show their joy? (By song, music with harp and zither, cymbals, dancing.) --Summary, Wine harvest.

\subsubsection{b) Geographical description of a specific place}

\emph{Nazareth}. The home of Jesus was a small, unimportant place, but beautifully situated in the midst of the province of Galilee. You can see this from the picture of the present-day Nazareth (En Nasira; Picture by Langl). A description of this picture shows the homes of the nearly 7,000 inhabitants scattered over the side of an almost unscalable mountain. Today the majority of the people are Christians. The place is surrounded by gardens of olive trees, fig trees, palms, and cypresses. The houses are of stone and have flat roofs, as is the custom in the East. (Reminder of David.) In the foreground we see a woman carrying a water jar upon her head. She is coming from the tree-shaded Virgin-spring which bubbles in the valley before the town. It is the only spring in Nazareth; and it is here that Mary, the mother of Jesus, may, with the child Jesus, have often drawn water. Farther down, but still on the heights around the town, are the orchards, with cultivated fields between them. Farther in the background are the less fertile heights which serve as pasture for the sheep and the goats. Beyond the height stretches the road to the fruitful Plain of Jezreel, watered by the Kishon (Reminder of Elijah), and surrounded in the blue distance by the ranges of the Carmel mountains. If we could climb to the summit of the mountain with those two Franciscan monks, whose
\quellseite{606}
cloister, however, is not visible on the picture, we should see to the far South the mountains of Gilboa. (Reminder of Saul’s death.) We should farther behold the steep heights of Tabor in the East, the mountains of Galilee to the North, and beyond them the peaks of Lebanon and Hermon, snow-covered even in summer, and the Mediterranean Sea in the far distant west.—Bittdorf, \emph{Methodik}.

\section{2. Instruction in the Catechism}

\subsection[1. A Catechesis to Begin Instruction]{1. A Catechesis with Which to Begin the Instruction of the Catechumens}

The catechist sings with the children the first three stanzas of the hymn, “\emph{Little Children Come to Jesus}”, reads Mark 10:13–16, and says a short prayer. Then he begins his catechization in the following manner:

\emph{Aim: Dear Children! We would hear today that the Saviour wants to have also you children with Him to bless you.}

\subsubsection{1. Developing the truth from the intuitional material} Among all the stories of the Bible I know one which should be especially dear to you and all children. It is so beautiful that the picture of it should hang in every Christian school. I know of none which shows more clearly that the Lord Jesus dearly loves children. Of what am I thinking? Yes, the story of Jesus blessing the little children. That is a true children’s story, and one can imagine nothing more lovely. We want to study it more closely. It was upon a day when many people had again followed Jesus and He had been ceaselessly teaching them and healing their sick. It was, therefore, a day of extreme exertion for the Lord. Finally came the mothers, pressing through the crowd and bringing their children to Jesus. The larger ones they led by the hand; the smaller ones they carried in their arms. What was the Lord Jesus to do for these little ones? He was to touch and bless them. That was a pious wish. Those mothers had seen that the sick were healed; that the blind received their sight, and that the lame walked when Jesus laid His hand upon them. They had also seen that it was good to be in the presence of Jesus, and had beheld how great a blessing came from Jesus upon all whom He touched. For whom did they wish such a blessing? For their children. Was there something lacking, or were they sick? No. Then what kind of blessing did these mothers want for their children if they were not sick? A blessing for their
\quellseite{607}
souls. Such it would have been, even if these mothers had not known what they wished for their children. They surely wished their children to become good, pious children. But Jesus had labored the whole day and was weary. What does one desire when he is weary? He seeks rest. What would Jesus then gladly have done? He would gladly have rested. Yes, we can well think so: for although our Saviour was true God, He was at the same time true man; and, while here on earth, He often became weary like other men. Who else must have thought that the Saviour wished to rest? The disciples. What did they therefore do when they saw these mothers coming with their children? They forbade them. Yes, they motioned to them with their hands that they should stay away. And as that did not keep them back, they spoke to them with hard words. And when that did not stop them, the disciples threatened them with force. What, under the circumstances, did it seem impossible for these mothers to do? To get near to Jesus. It seemed as though they would have to leave without accomplishing their purpose. Who, then, took the matter in hand? The Lord Jesus. How did He feel when He saw what was happening? He was much displeased. With whom? With the mothers who were coming to disturb His rest? Oh, no, with the disciples. Why? Because they would not let the mothers come with their children. What did Jesus say to His disciples? Suffer the little children to come unto Me, and forbid them not. He knew that these mothers were eager to bring their children to Him; and what must He have wished when He rebuked His disciples? He wanted to have the children with Him. But was He not weary in consequence of His preaching and healing of the sick? Certainly He was weary. If He was weary and yet wanted the children to be permitted to come to Him, which was to Him most important, to rest or to have the children with Him and to rejoice with them? To have them with Him and to rejoice with them. Yes, children, we have such a Saviour who forgets all weariness when souls come to Him and desire His blessing. Even when it is who that comes? Even when it is a child. Therefore we learn from this history: \emph{Jesus wants the children to be with Him, even when He is weary. He always has time for them.} Let us say that together. Again! Again! Anna, say it alone. Frank, repeat it. Now I shall write it on the blackboard.

The disciples did not wish to let the mothers come to Jesus with their children, because He was tired and needed rest. Can you give
\quellseite{608}
another reason why they probably wanted to turn these mothers away? Because they thought that Jesus could do nothing for these children. That He could be a Saviour for grown folk they had learned from His dealings with them and others, but how could He help these little ones, who could not even understand what He said? However, who must have trusted that He could give something to these children? The mothers must have so trusted. Otherwise what would they certainly not have done? They would not have tried so hard to force their way to Him with their children. Upon whose side did Jesus place Himself, that of the disciples or that of the mothers? Upon the side of the mothers. What did He say that He could and would give the children? The Kingdom of Heaven. What are His exact words? For of such is the Kingdom of Heaven. Not only grown folk, but also children can be His disciples and therewith members of the kingdom of heaven. No one—no one will be shut out unless he shuts himself out. The Lord Jesus has expressed Himself emphatically with regard to this, in His words to His disciples. Who knows what these words are? “Verily I say unto you, whosoever shall not receive the kingdom”, etc. Like whom must therefore the grown folks become before Jesus can use them in His kingdom? Like the children. The Lord Jesus did not thereby mean that the disciples and all adults must again become small and weak like the children, but He means that the grown-ups must become as open and receptive for the good and heavenly as children are in comparison with their elders before He can use them in His kingdom. The disciples must have opened their eyes wide at these words. They thought that Jesus could do nothing for whom? For the children. And who must have rejoiced at these words of Jesus? The mothers. Yes, indeed, for they perceived that Jesus could give the children much more than they thought. How did Jesus prove before all people that He can be something to the children? What did He do with them? He blessed them. More exactly! He took them up in His arms, put His hands upon them, and blessed them. There you see the true friend of children, who casts no one out, who also presses the smallest and least to His heart and would give them all the blessings of His heavenly kingdom. Let us note: \emph{Jesus can and will give also children His heavenly gifts.} Let us repeat that together a number of times. John, say it alone. Louis, do the same. And now I shall also write this upon the blackboard. How many things have we thus far learned from our history? Two. What is the first? Jesus
\quellseite{609}
wants the children to be with Him, even when He is weary. What is the second? Jesus can and will give also children His heavenly gifts.

\subsubsection{2. Comparison} The Lord Jesus always showed that He had a place in His heart for the children. Take your New Testaments and turn to the Gospel of Matthew, chapter 18. Read verses 1—4. About what were the disciples at strife? Which one of them was to be greatest in the Kingdom of Heaven. Where did they therefore all expect to be together? In the Kingdom of Heaven. They felt certain of that; but about what did they contend? Which of them was to be the greatest—the chief, in the Kingdom of Heaven? That one or the other of them might never get into heaven was not questioned by any of them. They had forsaken all and had followed Jesus, and how could they be shut out when Jesus should as king set up His kingdom? What did the Lord Jesus do to awaken the disciples from their dream of security? He took a child that happened to be in the house and placed it in the midst of the disciples. And what did He say to them? Read it in verses 3 and 4. “Verily I say unto you . . . . . . . greatest in the Kingdom of Heaven.” About which question should the disciples therefore not argue? Who of them should be greatest in the Kingdom of Heaven. But of what should they strive to be certain? That, above all things, they should enter the Kingdom of Heaven. And on what condition only were they to enter it? That they become converted and become as little children. As small and weak and ignorant as little children? No, as willing, simple, and humble as they. Who, therefore, will go into the Kingdom of God before them, according to Jesus’ words? The children. Hence, Jesus must have determined to receive also whom into His kingdom? The children. Certainly. What He here says is the same as when He says: Suffer the little children to come unto Me; for of such is the Kingdom of Heaven.—Later, when Jesus rode into Jerusalem on Palm Sunday, He declared that the children had understood Him better than many grown people, and that they gave Him great cause for joy. Open your Testaments at Matthew the 21st chapter. John, read verses 15 and 16. “And when the chief priests . . . . . perfected praise.” You know that when Jesus rode into Jerusalem on Palm Sunday a great procession of people praised Him. And that was so contagious that also who joined in? The children. What did they shout? “Hosanna to the Son of David.” They did right, for Jesus desired to set up His kingdom, in which they too should be welcomed and receive all that they needed. Who only
\quellseite{610}
stood by and grudged Jesus this joyous reception? The chief priests and scribes. What, according to their opinion, should Jesus have done? He should have silenced the children. But what did Jesus reply? “Have ye never read . . . . . . . perfected praise?” Yes, He took the children under His protection and said that they were doing precisely what they should do; and that what the Psalmist once said of the babes and sucklings, they were doing, namely, praising God for setting up His kingdom. Thus we again see how much Jesus loved the children and how He had determined to receive them into His kingdom and to bless them therein. That which we learned from the history of Jesus blessing the little children is entirely correct. What was it? 1. Jesus wants the children to be with Him, even when He is weary. 2. Jesus can and will give also children His heavenly gifts.

\subsubsection{3. Valuation} In what did Jesus desire to give children a part? In the Kingdom of Heaven. Is that anything so very great and precious? We will consider this more closely. Tell me, is it something great to be the child of a king? Yes, it is something great. The child of a king has everything it needs. There are many servants who wait upon it and fulfil every wish; and when the child has grown up, sooner or later, he will help rule the land. The greater the kingdom, the greater is the glory of such a king. When it was expected that a son would be born to Emperor Napoleon, they had a golden cradle ready, in which they afterward laid the child. Whose kingdom is greater than that of any king or emperor on earth? The Kingdom of God. Who will therefore have it better and more glorious, one who has part in an earthly kingdom or one who has part in the Kingdom of Heaven? He who has part in the Kingdom of Heaven. Where is the Kingdom of Heaven?—You can not tell? The Kingdom of Heaven is where Jesus Christ is. He who has Jesus Christ as his friend and helper is in the Kingdom of Heaven, and he is far better off than the son of the richest king upon earth. Jesus can bless, protect, and comfort, better than any king upon earth. What did Peter say when He saw Jesus upon the Mount of Transfiguration? It is good to be here. Peter was right, children! Nowhere is it better than with Jesus. Therefore we sing in church, “\emph{Who is, Jesus blessed, Like to Thee, sweet Rest?}” Ask the disciples; ask John, who lay upon His breast; ask the children whom He embraced and kissed; ask Mary the sister of Martha, who knew nothing sweeter than to be with Jesus! The person who is with Jesus is blessed; for he has all that he needs, even when he becomes sick and poor, even in trouble and under the
\quellseite{611}
cross. Instead of saying, The Kingdom of God is where Jesus is, we can also say, The Kingdom of Heaven is there where we have God as our Father and are become His children. When is a child happiest and most blessed,—when it is with whom? When it is with father and mother. Why were Adam and Eve so happy and blessed in Paradise? Because God was a father to them and they could speak to Him as children. When shall we be happiest? When God is our father and we are His children. Yes, children, that is the true blessedness. Even Heaven itself would be of no benefit to us, if God would not be there as our father and we as His children.

\subsubsection{4. Application} Of those children in the story whose mothers brought them to Jesus, He said: “Of such is the Kingdom of Heaven.” and we have just learned that to be in heaven means nothing less than to be with whom? With the Lord Jesus. And to have whom as father? To have God as father. And what is man's condition when he is with Jesus and has God as his father? He is blessed, then and then alone. Have you been brought to Jesus? Yes. When were you brought to Jesus? When I was baptized. Who brought you to Jesus? My parents. What must they have thought of Baptism if they had you baptized? They must have thought that Baptism was good and beneficial. Yes, they knew that, in Baptism, the Lord Jesus would bless you as He blessed those children who were brought to Him. They had learned from the words of Jesus that He wanted also whom with Him? The children. And that He could and would help also whom? The children. They had learned that He wanted to give the children a share in what? In the Kingdom of Heaven. They well knew that to be in the Kingdom of Heaven meant the same as what? As being with Jesus and having God as their father. And it was also known to them that you would be in what condition, too, if you were with Jesus and had God as your father? That we should be happy and blessed. Yes, that is why your parents brought you to Jesus in Holy Baptism. They wanted you to be with Jesus, to have God as your father, and thus to be happy and blessed children. See how well they meant it! But is it really true that, by Baptism, one gets God as his father and becomes His child? Yes. The Apostle Paul was strongly convinced of this truth; for when he writes in his letter to Titus (open New Testament at place): “According to His mercy He saved us by the washing of regeneration, etc.”, he means nothing else than this. Before he believed on Jesus and was baptized by Ananias in Damascus, he was not with Jesus, God was not his father, and he certainly
\quellseite{612}
was not blessed. He was very unhappy and had no peace with God. But when he believed and was baptized, he was as if new-born, happy, and blessed; for then he knew that God was his father and he was God’s child, and could say: Now, with all my sins and in every need I can go to Him with all assurance and confidence. Paul also thought of this when he wrote to the Galatians (find place in Testament): “Ye are all the children of God, etc.” All who believe in Christ and are baptized are God’s children, and Jesus surrounds them all with His presence and protection, as the clothing we put on covers the body. Paul’s baptism was considered by him to be the great turning-point in his life. Before that, he was without Jesus, without the Father, and therefore, instead of being blessed, lost and condemned; after being baptized, he was with Jesus and the Father, and therefore blessed. So, then, when did Paul’s blessedness begin? With his baptism. Not only when, therefore, did he consider himself blessed? After his death and in heaven. Hence he did not write: God shall save us, but; God hath saved us; He has already done so. And as it was with Paul, so, dear children, is it with us, with you. Also you are saved through Baptism and have received full salvation; for you have been brought to Jesus. He bestowed His kingdom upon you in that He became your Saviour and Friend, and in that God became your Father. The salvation which you shall have in heaven grows out of that salvation which was given you in Baptism, as the apple-tree grows from the seed. Whoever is not saved on earth will not be saved in heaven; for to be saved means nothing else than to be with Jesus, to have God as Father and to be His child. Let us repeat that! Again! By what means were you brought to Jesus, get God as your Father and become His children? Through Baptism. Oh, then Baptism is something great and glorious. A certain French king once said: “The three handfuls of water with which I was sprinkled in my baptism are more precious than the royal crown which I now wear upon my head.” A baptized child is like the child of a king in its cradle; salvation is just as certain to him as is the crown to the king’s son. Nay, he is much more fortunate, for the king’s child must wait many years after its birth before it receives the crown, while the child that is baptized receives salvation immediately in Baptism. The Lord Jesus takes it and blesses it, makes of it a child of God, which can always with all confidence go to God as a child to its Father. So you who are baptized are not the losers thereby. You have received in Baptism greater salvation than those children
\quellseite{613}
who were brought to Jesus by their mothers. And if you have been brought to Jesus by Baptism, have God as your Father and are in truth saved, what should, for that reason, be your greatest care? That we hold fast to this salvation and not lose it. That means that you remain what? That we remain God’s children. Yes, children, remain with Him, with Jesus, with God your father. And to this end may these hours of religious instruction help—that you remain the children of God!

Now open your catechisms and read, page 27: (“That . . . I may remain a child of God”) (the lesson to be assigned). Also read the three Scripture passages found in that connection. You can all understand this now very well. Memorize the Scripture passages for tomorrow. But we first want to repeat each of them several times together. Read carefully that which is printed in large type. I shall ask you about it tomorrow. And now, since we have been brought to Jesus by Baptism, have God as our Father, and therewith are in possession of salvation, let us sing a hymn on Baptism. I will read the first verses and then you may read them. (Here the catechist should paraphrase the difficult words.) And finally we will sing them together.

\subsection{2. Second Catechesis (outlined)}

\emph{Aim: The Holy Scriptures the Book of all Books.}

\subsubsection{1. Developing the truth from the intuitional material at hand:}

From the Scriptures Jesus shows the sad disciples on the way to Emmaus, who believed everything to be lost and who could nowhere find comfort and new hope, that they have no reason to mourn. He shows them the way God took to redeem mankind and make them His children, and also the ways they now must go. Their hearts rejoice and their eyes become bright. They now understand the ways of God.

\subsubsection{2. Comparison:} Philip also uses the Scriptures when he would show the eunuch of Ethiopia the same, and the latter is thereby led to Baptism and joyfully goes his way.—The Bereans search the Scriptures for themselves; and each new discovery makes them the more certain that Jesus is the Redeemer through whom alone they can become and remain God’s children. The grandmother of Timothy instructs him in the Scriptures from childhood, and Paul directs him to them when he has reached manhood, as the book that can make him wise unto salvation.

\quellseite{614}
\subsubsection{3. Valuation:} Therefore the Scriptures are greater and more precious than all other books. Therein we are instructed concerning the highest and most important things. Therein it is not mere men, but God who speaks to us. 2 Pet. 1:21. Ask the psalmist and the prophets (Scripture verses); ask Luther; ask many great ones in the world, and many a humble mother. You will find that they have all learned this: By means of the Bible God Himself speaks to us, instructs, comforts, admonishes as none other. If we would show the heathen the way they can become the children of God, we take the Bible in hand. If we ourselves want to learn this fact anew, we listen to the preaching from the Bible, and also search the Word ourselves. What do our hymns say?

\subsubsection{4. Application} And now we also want to learn from the Scriptures how we became God’s children and how we may remain such.—Search the Scriptures.—Show how that is done.—Compass—if we would go the right way, we must always—give individual cases—examine them.

\subsection{3. A Catechesis on the First Commandment}

The catechist sings with the class several stanzas of the hymn, “\emph{If thou but suffer God to guide thee}”, etc., reads Exodus 14, and says a brief prayer. Then he reviews the previous lesson, whereby the present one is introduced.

\emph{Aim: We would learn that children of God trust in God above all things.}

\subsubsection{1. Developing this truth from the intuitional material}

To what did God call Moses when He appeared to him on Mount Horeb and spoke to him out of the burning bush? Moses was to lead Israel out of Egypt. Yes, he should be their deliverer. That was a difficult task. Why was it a difficult task? Because Pharao did not want to let them go. And why was it for Moses especially a risky matter to return to Egypt and lead Israel out? Because he was once compelled to flee from Pharao. What did Moses say when God assigned him this difficult task? He did not want to go and raised many objections. What did he finally say to God? Send whom Thou wilt. Moses was not altogether wrong. Had he gone to Pharao in his own strength and attempted by his own wisdom to lead Israel out of Egypt, he had better not have undertaken it. He would have failed miserably. When did Moses truly experience this? When,
\quellseite{615}
forty years before, he had tried it by his own wisdom and strength. But now he wanted to do it under God’s commission. That was entirely different. What could he rightly expect when he was acting by God’s order? That God would help him. Certainly God also promised him that He would be with him. That was a glorious promise, and Moses could say to himself: Yes, if God will go with me, I will make the venture. He is mightier than Pharao and everything else. He is the creator of heaven and earth. He can show me how to begin and how to accomplish the work. I could have no better companion. If he is for me, who can be against me? When God said to Moses, I will go with thee, He after a manner stretched out His hand and said: I will lead thee. And what did Moses finally do? He laid his hand in God’s hand that He might lead him. And of what could Moses be certain when God was leading him? That he would be rightly led. Yes, Moses trusted God; for when one is given a hard task and he places his hand in God’s hand and trusts that He will lead him aright, that he will succeed with God’s help, we say of him that he trusts God. We will write that on the blackboard: \emph{He who trusts God, places his hand in God’s hand that He may lead him, and is certain that with God’s help he can accomplish the most difficult task.}

As Moses went to Egypt, which words of the Lord must have been constantly ringing in his ears? The words: I will be with thee. Yes, they were his comfort, and always gave him fresh courage. When his brother Aaron met him on the way he was very glad; but upon whom did he rely more than upon Aaron? Upon God the Lord. When would Moses later especially have relied upon God’s assistance? When he had to go to the children of Israel and tell them that God had appeared to him and commissioned him to lead them out of Egypt. But on what occasion to an even greater degree? When he had to appear before Pharao. Yes, all kinds of thoughts must then have arisen in Moses’ mind, such as these: Suppose he recognizes you, and still remembers why you fled before, and imprisons you; or has you killed! With such thoughts as these it would not have been surprising if he had lost courage and turned away. Why did he not do so? Why did he go calm and unafraid to Pharao to fulfil his mission? Because he trusted in God. He heard God saying again and again, I will be with thee, and he gripped His hand the tighter. And what did he trust God to do? To help him and lead him aright. At any rate he was sure that as long as God stood by him he could
\quellseite{616}
not fail. He must have said to himself, “Commit thy way unto the Lord; trust also in Him, and He shall bring it to pass.” Did God stand by him and champion his cause? Yes. True, in the beginning it did not seem so. What reply did Pharao give when Moses demanded of him that he let Israel go? I will not let Israel go. And what did Pharao say when Moses called upon God the Lord? Who is the Lord that I should obey his voice? I know nothing of the Lord. And when Pharao refused again and again, and finally forbade Moses to appear before him, it must have been difficult for him to trust God. All seemed to be in vain. Did he then give up his trust in God? No. What rather did he do? He prayed to God. He kept his hand in God’s and continued to hope that in the end God would give him success in his undertaking and lead Israel out. Even though he might have to wait long and Pharao might refuse again and again, he said, “Nevertheless I am continually with Thee: Thou hast holden me by my right hand. Thou shalt guide me with Thy counsel, and afterward receive me to glory.” With such thoughts as these he strengthened his courage and assured himself that since God had bidden him take this way, He would see to it that he did not fail. We note: To trust God is to put one’s hand in God’s hand, to leave it there and to hold it ever fast, even though one hindrance come upon the other and make the performance of our duty more and more difficult. (Blackboard).

After Pharao had let the Children of Israel go, where was Moses trust in God again put to the test? At the Red Sea. Other trials of his trust had not been wanting, but the hardest came at the Red Sea. When God showed him the road he must travel, he had to trust Him. How did God show the road? He sent a pillar of cloud before the Children of Israel. Did He in this way guide them to Canaan by the nearest way? No. At first it seemed as though God would lead them straight to Canaan, but in which direction did He suddenly guide them? Toward Egypt. Correct. He made a sharp turn southward and so led them farther and farther into the country of the enemy. To whom especially must strange thoughts have come? To Moses. It was probably hard for him to understand why they were being led that way, for Moses knew the way to Canaan well. What would he have done if he had followed his own understanding? He would have said, we are being led the wrong way, and he would have let go of God’s hand. What would he then, of course, not have done? He would have no longer trusted God. What does it mean to trust God? To
\quellseite{617}
place your hand in God’s hand and let Him lead you. Yes, and to be sure of what? That His way is the best. But Moses did not cease trusting God. He could not understand why God was leading them this way; to his mind it was all wrong. But he leaned more upon God than upon his own understanding. And what must Moses, more than ever before, have said to himself? Think of the Scripture passages you have already learned. Commit thy way unto the Lord; trust also in Him and He shall bring it to pass. And if he had been acquainted with our hymnal and learned our hymns, which one would he have thought of? “\emph{Commit thy way confiding, When trials here arise.}” And when the morning revealed the fact that they had travelled southward the whole night and were going farther into the land of the enemy, and his reason was telling him more and more plainly that this was wrong,—what did he then do with God’s hand? He still held fast to it. Correct. He clasped it so much the tighter and said, “Nevertheless, nevertheless, I am continually with Thee”, etc. But you said before that Moses’ trust was given its severest test at the Red Sea. Whither did God finally lead the Israelites? To the Red Sea. Yes, the Red Sea lay before them, and what appeared to the right and left? High mountains. And who came behind them? Pharao and his army. Israel’s case seemed desperate, shut in on all sides like a mouse in a trap or a fish in a net. There seemed to be no way of escape. And who was it that had thus led Israel? God himself. And had God led them wrongly and brought them into misfortune? It appeared so. Who began to believe this? Many of the people. What did they say to Moses? “Why did you lead us out of Egypt? Were there not graves in Egypt, that we should perish in the wilderness?” But who could not believe that God had led them wrongly? Moses. He had beheld too many of God’s wonderful ways, and could not now cast his trust in Him away. It is true that his spirit was deeply grieved and that he could find no way out of the difficulty, but he did not give up hoping in God. God had called him to deliver the people of Israel out of Egypt; He had led Him in the past, and he could not believe that God would now permit him and his people to perish. He held fast to the belief that God knew the way out of the trouble. And when his soul trembled and feared, he said to it: “Why art thou cast down, O my soul? and why art thou disquieted within me?” etc. Here we first clearly see what true trust in God is. He who in truth trusts God, is also
\quellseite{618}
certain that God’s way is the right one even when? Even when it leads into deepest distress and one can find no way of escape. Repeat this a number of times. \emph{Who trusts God is even then entirely sure that God’s way is the best even then when it leads into deepest distress and one can find no way of escape.} I will also write it upon the blackboard. Moses had strong trust in God. What did he do in the strength of such confidence? He cried unto God. What did he say? Lord, help us! Certainly, he reminded God of His omnipotence by which He can save us out of every trouble, and of which promise did he undoubtedly remind Him? That He would be with him. He reminded God of His faithfulness, according to which He must now keep His word. “I have trusted in Thee; Thou canst not and wilt not let me be put to shame”—thus he must have called to God. And what was God’s reply? “Lift thou up thy rod, and stretch out thine hand over the sea, and divide it; and the Children of Israel shall go on dry ground through the midst of the sea. And the Egyptians shall know that I am the Lord.” Oh, that was a new trial of his trust. He, a weak human being was to divide the waters of the sea and make a way through for his people. Did his faith stand this test? Yes. What did Moses do? He stretched out his hand over the sea. Behold him standing there by the side of the sea, with his hand holding the rod stretched out over the waters. His reason told him: Moses, you are making yourself ridiculous; never before has a man parted the waters of the sea! But Moses stood firm in this confidence: \emph{He trusted that God could even work a miracle to save His people.} That was the highest point of his trust. (Blackboard.) And see how contagious it was! What did the people do when they saw the parted waters standing like walls on either side and the way open through the midst of the sea? They went through. What trust in God! They, like Moses, laid their hands in God’s, that He might lead them. And what did they trust? That God’s leading would be the best even though it led through the midst of the sea. Children, that is trusting in God above all things, when we, like Moses, first do what? When we lay our hand in God’s hand, that He may lead us and are sure that with God we can perform the most difficult task. In the second place, when we do what? When, like Moses, we leave our hand in God’s hand and only hold the tighter when one hindrance comes upon another to make our task more difficult. Thirdly, when we do what? When we are certain that
\quellseite{619}
God is the best guide, even though we come into deepest distress and can find no way out. Yes, and what is now the last thing this trust in God clings to? That God can work a miracle if He will, to save His flock.—It was not strange that God was pleased with Moses; for he trusted God as but few trust Him.

\subsubsection{2. Comparison}

We see, however, that still other men and women trusted God. Name a king of Israel who trusted God above all things. David. When did David show his trust in God? When he went forth against Goliath the giant. What did he, accordingly, say to Goliath? “You come out to me with sword and spear, but I come in the name of the God of Israel.” Here also there was a great task to perform. In order to perform it, David laid his hand in God’s and trusted that He would help him slay Goliath and so fulfil his duty. When did David trust God in suffering? When Saul persecuted him for years. There he laid his hand in God’s hand and was sure that God would do what? That He would lead him aright. Yes, and even though he was forced to wait for years and years before obtaining the throne which had been promised him, and was forced to flee from place to place, what did he still firmly believe? That God was leading him aright. What, therefore, did he sing in the twenty-third Psalm? “He leadeth me in the paths of righteousness.” When did Abraham show his trust? When he left his homeland. Yes, when he left his own land to go into a land the name of which he did not even know. Then he laid his hand in God’s hand, that He might lead him. And of what was he certain? That God would be the best guide for him. Of what hymn stanza are we very strongly reminded when we see Abraham thus going forth into the strange, unknown land? “\emph{Commit thy way confiding}”, etc. What did Abraham thereby show? His trust in God. Here again we see that it dealt with the performance of a great task, which Abraham undertook, upheld by true faith in God. Which one of Abraham’s descendants also trusted God through long, hard trials? Joseph. When did his trials begin? When he was sold into Egypt by his brethren. And when did they end? When he became ruler over Egypt. How many years was that? More than 13 years. Oh, that was a long time of trial. What thoughts must certainly have often come to him? That God had forgotten him. When especially did such thoughts come to him? When Potiphar had him cast into prison. Again when? When the cup-bearer was freed and forgot his promise. But what did Joseph always do with such thoughts? He always overcame them. Yes, he trusted God. Even in the greatest trials he did not let go of God’s hand; but what did he do? He held it all the
\quellseite{620}
tighter. What confidence in God did he thereby show? That He had not forgotten him, appearances notwithstanding, and that He was the best guide nevertheless. Thoughts such as Paul expresses in the passage, “We know that all things work together for good to them that love God”, must often have passed through his soul. And thus he showed that the true children of God should not throw away their trust even when? Even when the distress is great and long continued. When we now see that Abraham, Joseph, Moses, David trusted in God above all things—and we could mention many others, especially our dear Dr. Martin Luther—who should then trust God? All who would be God’s children. Yes, to be God’s child and to calmly and trustfully place one’s hand in God’s that He may lead one, and to be sure that—even in the greatest trials—He is the best guide: these things necessarily belong together. May God’s children then trust no one but God? Oh, yes. In whom does the child trust? Its parents. In whom the pupil? His teacher. In whom the king? His people. In whom the sick? The physician. But the true children of God, whom do they trust more than any one else? God. Yes, they trust in God above all things. They regard Him as their only refuge and help. They are certain that others can help them only when who permits? When God permits. Hence, what kind of father do we call God in the First Article, because we Christians trust in Him that He will always lead us aright and help us in every need? The Father Almighty.\footnote{Here it would be better to stop and postpone the Valuation and Application until the next time. Such important fundamental truths as are here dealt with must be thoroughly considered, which may well take half a week. If the catechist has time, he may thoroughly go through the hymns “Commit thy way confinding” and “If thou but suffer God to guide thee”. Much else in the explanation of the Commandments can be briefly treated—the negative more briefly than the positive.}

\subsubsection{3. Valuation} Did Moses receive any benefit from his trust in God? Yes. Who went with him when he trusted in God’s help and went to Egypt? God went with him. How did that show itself? By the signs and wonders which he did before Pharao. By whose power had he forty years before tried to free the Children of Israel? By his own power. What did he then accomplish? Nothing. What was he then forced to do? He was forced to flee. But what was he now able to do, when, trusting God, he undertook the great work? He could lead Israel forth. What effect did his trust in God have upon him? It made him strong. Yes, trust in God makes one strong and
\quellseite{621}
courageous. We see this also in David. How large was David, compared with Goliath? He was small and weak and wholly without experience in handling weapons of war. He was so weak that who thought it impossible for him to conquer Goliath? His brethren and King Saul. But with what success did he nevertheless meet? He killed the giant Goliath. How did his trust in God make him? Strong and courageous. And fearless, we may add. Fearlessly Moses went before Pharaoh, fearlessly David went against Goliath, fearlessly Abraham went forth into the unknown country. When, therefore, can one undertake a task with courage and strength, and without fear? When one trusts God above all things. (Of course, there is one thing we must not forget. Who commanded Abraham to go into that strange land? God. Who sent Moses to Egypt? The Lord God. Who bade David to go forth against Goliath? God told him, for He sent him to the camp at the time when he had to hear how Goliath blasphemed God. By whom, then, must the task be imposed if one would assume it in God’s name and be sure that God will make one strong, courageous and fearless? By God. One has no right to push oneself where one is not called and then expect God to help him. To thoughtlessly thrust one’s life into danger and then expect God to help, is not trusting God, but tempting Him, and that is sin.) Upon what then does all depend when God gives us a difficult duty to perform? Upon this, that we undertake it trusting in God. When Lincoln sounded the call for troops after the shelling of Fort Sumter, he had to contend not only against the war-like Confederates but also against a strong opposition in his own country. Well might he have trembled. But fearless, strong, and courageous he went to war; for it was not he that had provoked the war. He could trust in God. His good conscience and trust in God made him courageous, fearless, and strong. Now let us repeat together: \emph{Trust in God makes us fearless, courageous and strong.} We will say it again. (Blackboard.) Of this the prophet once spoke, who himself always depended upon God for all things. It was the prophet Isaiah. Read Isaiah 40:31. Anna, read the verse, “But they that wait on the Lord”, etc. O children, all good people who trusted in the Lord have experienced how they received new strength. More than once they became so weary that it seemed as though they could not go a step farther; but what did God give those who trusted in Him? New strength. What could they therefore do? They could mount up like an eagle. How does the eagle fly? With power and to a great height. Yes, he seems to be absolutely tireless. And what else can they do? Run without becoming weary, walk without becoming
\quellseite{622}
faint. We see this in Elijah when he ate the food which God gave him; in Moses, who forgot all his weariness when he heard that he should stretch out the rod in his hand and lead Israel through the sea dry and safe. Who would not gladly trust God when he thereby becomes fearless, courageous and strong?

If we look at Joseph, we shall learn another truth about the blessing of trusting God. How long did Joseph have to suffer trials? More than 13 years. That is a long time, as long, children, as you have lived. And when things go wrong, time passes, oh! so slowly. A single day seems very long, and how long would 13 years seem. Think of a sick person who must lie still for 13 years or more! Joseph might easily have become discouraged if he had not done what? If he had not trusted to God that He would eventually help him and restore him to honor. And because he trusted God, he would think like that other pious man who was also compelled to endure much, but who repeatedly said to himself, “I will trust in Thee”, etc. What did Joseph learn because he was sure that the end would be good? He learned to wait until God’s hour came. And that was no impatient waiting. When it has rained for a week, we often can not wait until the sun shines again. What kind of waiting is that? That is an impatient waiting. Joseph waited patiently and calmly. What do we generally do when we wait impatiently? We grumble and complain. Yes, sometimes our fault-finding and complaining is very much in evidence. On the other hand, what was Joseph like, since he did not murmur or complain? He was quiet and calm. And why could he be so quiet? Because he trusted in God above all things. Find the 62nd Psalm. John, read the first verse: “Truly, my soul waiteth”, etc. Who also could speak like the Psalmist? Joseph. His soul waited upon the Lord, of whom he was certain that He would help him when the proper time was at hand. What, then, does trust in God tend to make us also? It makes us patient and calm. Correct. And it does so because it makes us cheerful. Only that person who is cheerfully confident, who without doubt and fear looks forward to the end, can wait, be patient and still. Let us say that together: \emph{Trusting in God makes us cheerful, patient, and calm.} (Blackboard.) Children, that is one of the greatest virtues, to be cheerful, patient, and calm in trial and suffering. To undertake a difficult task in a strong, courageous, and fearless spirit means much; but to remain cheerful, patient, and calm amid long drawn-out trials—that is much more. And in whom only shall you find both? Only in him who trusts God above all things. What, then, must we learn? True trust in God.
\quellseite{623}
What did God do with Abraham, who trusted in Him? He led him into the Promised Land. What did He do with Joseph? He finally exalted him and made him ruler over all Egypt. What success did Moses have because he trusted in God? He let him lead Israel out of Egypt. What did He let David do? He let him conquer Goliath. Therefore, what did He always do? He always helped. Yes, He has always done all things well, and brought things to a successful end, for whom? For those who trusted Him above all things. Provided they always did what? Provided that they laid their hand in God’s hand that He might lead them, and provided they were sure that He is always the best guide. Trust in God, accordingly, makes us not only fearless, courageous, and strong when we have to undertake a hard task; likewise, not only cheerful, patient, and calm when we must suffer long and much; but what else does it do? It leads to a successful end. To what does it lead when we have a difficult task? That we may accomplish it. From what does it keep us under long, hard trial? From murmuring and complaining. Yes, from murmuring and complaining against God and man; from doubt of God’s help. And to what good end does it lead? That we shall be finally delivered out of all trials and once more have happiness and joy, health and success. Let us say that together! Trust in God always leads to a good end. Always? Children, is that not saying too much? Think of Lazarus! He certainly trusted God, and how did matters turn out with him? Badly. How badly? His body was covered with sores. Was there any other misfortune? He was poor and had nothing to eat, but the crumbs from the rich man’s table. He was poor and sick; and, children, we should not forget this: he had to see others about him strong and healthy and that one before whose door he lay living in mirth and splendor every day. That made his sickness and poverty all the more bitter for him. You say that Lazarus trusted God, and that trusting God always leads to a successful end. Did it lead to a successful end in the case of Lazarus? Yes. But how long did he remain sick and poor? Until his death. Did he then come to a successful end in this life? Do you call that a successful end when one must remain in sickness and poverty until his death, while others are healthy and rich? No. How can you then say that the trusting of Lazarus led to a successful end? Because this life does not end all. What further happened to Lazarus? His soul was carried up into Abraham’s bosom. What was its condition there? Unspeakable happiness. Yes, God so highly honored him and surrounded him with so much happiness that he no longer thought of what? Of
\quellseite{624}
his sickness and poverty on earth. But only when did matters come to a successful end? Not until after his death. Then is it correct that trusting in God leads to a successful end? Yes, it is correct. But how should we not understand it? As if God already in this world always lets things turn out well. Oh, no, children. God often does as in the cases of Abraham, Joseph, Moses, and David. When does He often permit the one who trusts Him to enjoy success? Here in this life. And what does He often do with trials if one trusts Him? He brings them to an end and again bestows health or riches or honor. But He does not always do this. Many a one who trusted in Him was not permitted to see the success of his endeavors; more than one pious Christian has died in poverty and wretchedness. But when would He always help him, give him honor and bring things to a successful end? After death. Where does He take those who in this life have trusted Him above all things? To Himself, into everlasting joy and glory. We must therefore add something to the sentence “trusting God leads to a good end”, in order that it may not be misunderstood. What words would you add? In this life or in the life to come. Yes, we will so note it and repeat it together a number of times: Trusting God leads to a good end, if not in this life, then in the life to come! (Blackboard.) When, especially, will the true children of God have to prove their trust in Him? When God leaves them in distress until the end of their lives. For what especially will they then trust God? That after death He will bring them to honor and lead them to joy and glory. In what words will they thereupon take new comfort? Turn to Psalm 73. Frank, read verse 23. Yes, so they will say, now all the more. And if their spirits become disquieted, they will say to themselves as we read in Psalm 42:11. What do you read there, John? “Why art thou cast down”, etc. Or with what the Lord says in Isaiah 49:15. What is said there? “Can a woman forget her suckling child”, etc. (Here have the hymns read, “Commit thy way confiding”, and “If thou but suffer God to guide thee”, discuss them briefly. Also Psalms 23 and 91.) Yes, children, so great and important is trust in God. What does it do in the first place? In the second? In the third? Should you not also trust God above all things, place your hands in His that He may lead you, and trust Him that He is always the best guide? Is His guidance good even when the task He has assigned you is great and hard; when the trials He sends last long; even when He lets you die in wretchedness and poverty? Yes, even then.

\quellseite{625}
\subsubsection{4. Application} Such great tests of trust in God as we find in the lives of Moses, Abraham, Joseph, or David, are not demanded of you children. You are not required to forsake friends and acquaintances and go alone into a strange land, the name of which you do not know, as did Abraham. You will not be sold into slavery like Joseph and, though innocent, have to lie for years in prison like he; you are not required to deliver a people out of Egypt and to lead them through the sea like Moses; nor need you, like David, go out to fight a giant. Though in later days hard tasks and severe trials may confront you, they are as yet small and insignificant. But, nevertheless, what shall you soon have to show? True trust in God. If you have a difficult problem to solve or an examination is before you, and you are anxious and fearful, what should you do? Trust in God and hope in His help. Here, certainly, also another thing is necessary. What do you first have to do, if you know that an examination will shortly be held? Study diligently and work hard. Of course, for God has no pleasure in lazy people; He will rather let them fail dismally, so that they may finally learn to be diligent. But often, though you had been diligent, not only occasionally but regularly, from day to day, yet, how did you feel when the examination approached? You felt afraid. Yes, your heart beat and leaped violently. How, then, can you quietly and serenely go to the examination? If you say to God: Dear Lord, be with me and help me to give the correct answers. For what does trust in God make one? It makes one fearless, courageous, and strong. A youth was working for a farmer and was sent into the field with the horses and machine. He was only sixteen years old and had just recently come from Germany. He had but little experience with such young, mettlesome horses, and had seen the machine in operation but once. How would that work out? His heart was heavy as he rode out to the field. The young horses were worse than usual, for the flies were plaguing them. And when he tried to set the machine, he could not, for he had forgotten how it was done. He tried again and again, but without success; the horses refused to stand any longer, and he could not discover how to operate the machine. He wished he were back again at the house or with his parents in Germany. Then he happened to remember that in the First Commandment it says that we should trust God above all things. Do you know what he did? Reins in hand, he knelt down beside the machine and prayed: “Lord, help me! Thou hast often helped my mother when she called upon Thee. Help me also!” Then his heart grew calm and
\quellseite{626}
his head clear. The next time he tried it he succeeded, and before the day was over he had accomplished considerable work. See, boys, that is the blessing of trusting God. You will soon be confirmed and more than one of you will be going forth from your homes to the service of others or to another school. That may not be pleasant and your hearts may be troubled and heavy. But what did Abraham do when he left for the strange land? He trusted God. And how did he then feel? He became fearless, courageous, and cheerful. What can you do that you likewise may be of good courage? I can trust God. Yes, you can place your hand in God’s hand that He may lead you, and you can be certain that He will lead you aright. Does that mean that it will always go well with you? No. Often people will not be satisfied with you, however much you may strive to please them. Of whom should you then think? Of Joseph. Yes, he did his best and was sent to prison. But what did he not do, despite all this? He did not forget God. Still more. What did he do with God’s hand when things went wrong with him? He held it the closer. And of what was he confident? That God was the best guide. And that it would lead to what kind of an end? To a successful end. So you must think when you have done your duty and men are not satisfied with you. Or perhaps your parents become ill and die. Upon whom will you then rely? Upon your wisdom and cleverness, upon your strong, healthy body, upon your friends and acquaintances, or upon God your heavenly Father? Upon my heavenly Father. Why not upon your own wisdom and strength? Because they will not suffice. Why not upon friends and acquaintances? They may die, and even though they live, they can not always help. Who can always help, however great the need? Our heavenly Father. What kind of father is He? An almighty father. If He should let you become ill and even die, while the others about you are healthy and happy, what would be difficult for you? To trust in God. Certainly. But if you pray to Him He will preserve your trust in Him, and you will hold firmly that even then He is guiding you aright. At death pious children of God more than ever lay their hands in God’s and say: “Abide with me; fast falls the eventide”, etc. Even in the valley of death, accordingly, they are not alone; but who is with them? God is with them. What words of David’s Shepherd Psalm may they, accordingly, apply to themselves? “Yea, though I walk through the valley of the shadow of death”, etc. Why is, also then, God
\quellseite{627}
the best guide? Because the time to enter heaven has come. Yes, children, that is truly the best guide who leads us to heaven. If God then leads us so well, what should be and shall be our chief care? We will then trust in God above all things.—By all means, children, we will do that!—Open your catechism, page 31. There read what stands under the heading “trust in Him”; and again on page 32. Also read the Scripture passages plainly and with proper emphasis: Frank, No. 15, Anna, No. 16, Emil, No. 17, etc. Now, study the lesson carefully for tomorrow. (Passages which were quoted in the first part of the catechization have already been learned between the First and the Second Parts; Cf. page 598, note.)

\subsection[4. The Nature of Faith]{4. Outline of a Catechesis on the Nature of Faith}

Aim: We want to learn what it means to believe.

\subsubsection{1. Developing the nature of faith from the intuitional material:} From the story of the departure of Abraham from his fatherland the class is to be shown how faith has to deal with the invisible (unknown land), but only with such invisible things of which God has spoken to us (God spoke with Abraham about the unknown land); that faith clings to these as real and firmly relies upon them—so firmly, indeed, that it will give up everything visible and material for them.—

\subsubsection{2. Comparison} The martyrs, who gave up the earthly, tangible, perceptible, in order to gain the invisible, everlasting life. Heb. 11.—

\subsubsection{3. Valuation} Proceeding from Heb. 11, especially the life of Abraham, and terminating in Jer. 5:3; Heb. 11:6; Gen. 15:6; Hab. 2:4.—

\subsubsection{4. Application} Experiences from the life of youth.

\subsection[5. The First Part of the Third Article]{5. Outline of a Catechesis on the First Part of the Third Article}

(For the Juniors and Seniors of the Academic Department of a College.)

Aim: We cannot by our own reason or strength come to faith in the Lord Jesus. \subsubsection{1. Paul supplies the intuitional basis} The catechist, with the co-operation of the class, develops a sketch of Paul’s life up to the time of his journey to Damascus in a manner in keeping with the age and knowledge of those in the grades mentioned. He leads them to Tarsus into the house of Paul’s father, who himself had been a Pharisee; gives them an idea of the instruction on the Law which Paul here received; has them accompany Paul,
\quellseite{628}
the youth, to Jerusalem and sit beside him at the feet of Gamaliel; shows them by particular examples his zeal in the fulfilling of the Law (Phil. 3:1–11 to be read in this connection), making it clear to them, on the basis of Gal. 2:16, how only one thought filled the mind of Paul—his seeking after righteousness. He then presents comprehensively how knowledge and zeal were not wanting; how Paul outstripped the majority; how reason and strength were there, and yet total inability to know Christ in spite of them and to attain righteousness by his natural strength. No wonder he wrote 1 Cor. 2:14 and 1 Cor. 12:3. If Christ Himself had not intervened, and had not called to him on the way to Damascus; if He had not enlightened his mind and spirit while his eyes were blinded; if He had not revealed Himself to him (Gal. 2:11, 17), i.e., removed the scales so that he could see and know Jesus; if He had not let him be baptized by Ananias and thereby become justified and sanctified, he would never have come to know Christ as his Lord, and thus, in spite of all his knowledge and strength, remained a poor, lost soul. But when he arose out of the washing of regeneration he could confess: “I believe that I cannot by my own reason or strength believe in Jesus Christ, my Lord . . . . . . . sanctified and kept me in the true faith.” Now he could write Eph. 2:1 and Titus 3:5 and Rom. 3:28.—

\emph{2. For Comparison} the scene at Caesarea Philippi may be used, where the Master asks the question: “Whom do men say that I the Son of Man am?” and, upon the confession of Peter, rejoins: “Flesh and blood hath not revealed it unto thee, but my father which is in heaven”; the Bible reading and lack of understanding on the part of the eunuch of Ethiopia until Philip opens the Scriptures to him. One may also remind the pupils of Athens, which, with all its wisdom, could not find God, built an altar to the unknown God, and regarded the message of the risen Christ as ridiculous.—

\subsubsection{3. Valuation} Is it a piece of good fortune to be led to Jesus, to be brought to faith in Him? Ask the Greeks who desired to see Jesus; ask the converted Paul; ask the baptized eunuch who went his way with joy; ask Stephen, face to face with death by stoning; go out into the mission field and ask the converted heathen; listen to the hymns of Jesus which the Church sings!—

\subsubsection{4. Application} To know Latin and Greek, and to have other knowledge is good, but to know Jesus Christ is better than all knowledge.
\quellseite{629}
All wisdom and power do not lead to God; thank God that you have been baptized, thereby sanctified and brought to Christ. Are you still with Him, and can you say with Luther, “and preserved me”?

\subsection[6. The Conclusion of the Third Article]{6. Outline of a Catechesis on the Conclusion of the Third Article\footnote{From Berndt, \emph{Methodik des Unterrichts in der evang. Religion}, 1909, page 105 f. According to the analytical method, to which Luther’s explanation is the sum of the entire catechization, p. 523 f.}}

\emph{I believe in the resurrection of the body and the life everlasting.}

\subsubsection{1. With what Christians may comfort themselves at the grave}

a) How did Jesus comfort sorrowing humanity? The son of the widow of Nain, the raising of Lazarus.—b) How did the Lord confirm to them His comfort? “I am the resurrection and the life”, etc., certainty of the eternal life of Jesus; His resurrection the pledge of ours. He who is in Jesus has passed through death unto life.—c) How has God written this comfort already in the work of nature? Autumn, Spring, 1 Cor. 15.—d) How did pious Christians comfort themselves? 1 Cor. 15.—Paul (2 Cor. 5)—Ignatius of Antioch and others—The church hymns; “Jesus my Redeemer lives”, etc.;—Luther at the deathbed of his Lena.—e) How did D. Martin Luther proclaim this comfort? “In which Christian Church He . . . . . . at the last day will raise up me and all the dead.”—Resurrection of the body = all people.

\subsubsection{2. In what Christians may rejoice in the hour of death}

a) What does the Lord say about our joy in heaven? Parables.—b) Why is eternal life a great joy for the Christian? Rev. 21:1–7; Communion with Christ, God, all the redeemed, freedom from all inner imperfection and outer restraint.—c) How pious Christians would rejoice in eternal life. Stephen—Paul in his letter to the Philippians—Augustine and Monica (\emph{Confess.} XI, 10)—“Jerusalem thou City Fair”—and like hymns.—d) How does Luther speak of this joy? “And give unto me and all believers in Christ everlasting life.”

\subsubsection{3. How Christians conduct themselves on their way to heaven}

a) The earnest sermon preached to us on this subject by Christ. Parable of the Wise and Foolish Virgins, and others.—b) How the faithful disciples followed His admonition. Peter (Legend by Kinkel)—Paul.—c) How Christians proceed to heaven. “Blessed are they that mourn, for they shall be comforted.”—Appropriate hymns.

\quellseite{630}
\subsection{7. Questions for Review and Drill}

What can God demand of us since He has become through Baptism our God? When do we hold Him to be our God?—We should fear Him, i.e., always hold Him before our eyes as whom? We should love Him as who loved Him?—And we shall trust Him as who trusted Him?—When did Moses trust God above all things? When again? And especially when?—What did Moses do when he was commissioned to go to Egypt to save Israel? (He placed his hand in God’s that He might lead him.) What trust did he thereby repose in God?—And what did Moses not immediately do with God’s hand when Pharao would not let Israel go?—But what did he do all the more?—What did he further trust Him to do?—When did God seem to be an erring guide?—What did Moses then do instead of casting aside his trust in God?—How did he speak to his disquieted soul? Recall the Scripture passages you have learned! Frank, you, too, may repeat them.—By what can we see that Moses was right when he refused to cast aside his trust in God?—Was Moses ever sorry that he had laid his hand in God’s that God might lead him?—Show how David trusted God.—Show how Abraham trusted God.—From whose life especially have we seen that we should trust God even amid trials?—How long did Joseph’s humiliation last?—How did he comfort himself when it seemed as though God had forgotten him?—In what passage of Scripture has God Himself said that He will not forget His children?—What did Joseph say to himself when all evidence seemed lacking that God was thinking of him? Think again of your Scripture passages!—When, therefore, should we not let go of God’s hand, but hold it the closer?—And what trust shall we have in Him even in the longest trials?—For these two things belong to true trust in God: 1. Place our hand in God’s hand that He may lead us; 2. Be confident that He is the best guide, even when we have a hard task to perform or when long and severe trials rest upon us. In which Scripture passage are we admonished to commit all our ways unto Him and to trust in Him?—Which hymn also contains the same admonition?—We should trust God above all things also because a great blessing accompanies such trust. How did Moses practice this?—What did David do when he went out against Goliath?—Or Abraham, when he was told to go out into the unknown land?—In which Scripture passage are we told that trust in God makes one strong, courageous, and fearless?—Peter, you also may repeat
\quellseite{631}
it. Anna, likewise!—What effect did Joseph’s trust in God have upon him in his long trial?—But I know of another thing which should make it easy for us to trust in God. Who can tell it to me?—But does trust in God always lead to a successful end?—Even if we must die in poverty and misery?—From whose life have we seen that one may be pious and trust God and still have to suffer in wretchedness until death?—But why can we say that Lazarus’ trust in God led to a successful end?—What should we therefore add to the sentence: ‘Trust in God leads to a successful end’?—What must we truly do if trust in God is to make us fearless, strong, and courageous to fulfil our task; if it makes us patient and cheerful in trial, and if it always, in this or at least in the life to come, leads to a successful end?—In which psalm has David declared that God has always led him aright?—In which psalm would Moses help us to find joy in trusting God?—Which two hymns especially are calls upon us to trust in God? (When they have been learned and explained, other questions regarding them are in order). Show me how also you can exhibit true trust in God.—Louis, you also may show this.—Emil, likewise.—Yes, children, to trust God in all things is one of the main points in Christianity. It makes us courageous in life, cheerful, calm, and patient in trial, confident in death. May God Himself awaken this within our hearts and preserve it until our end!

\subsection{8. An Excursus Through the Catechism}

Thanking and serving. According to L. Schultze, \emph{Katechet. Bausteine}.\footnote{There is much to be improved before it is really fit to be used.}

It is a good thing to give thanks unto the Lord, dear children! Therefore we find so much about thanksgiving in our Catechism. Where do we find the word “thank” for the first time? In the First Chief Part, in the Second Commandment: Call upon His name . . . and worship Him with . . . thanksgiving.—For what shall we thank God? For everything.—Where is that written? In the First Article of the Christian Creed. In which words? “For all of which I am in duty bound to thank and praise . . . obey him.” Tell what that “all” comprises which is here mentioned. That God has made me and all my goods.—How can that “all” with which God preserves body and life, be expressed in two words?—Daily bread.—Who has taught us so to call it? Jesus in the Lord’s Prayer.—In which
\quellseite{632}
petition? In the fourth.—What do we therefore, according to this petition, owe Him also for our daily bread? That we receive our daily bread with thanksgiving.—We have already had three Chief Parts in which the giving of thanks is dealt with. Why is thanksgiving nevertheless so often lacking? Give me the entire explanation of the Fourth Petition. “God gives daily bread . . . to know and to receive our daily bread with thanksgiving.” What, according to this, should precede the giving of thanks? “That He would lead us to know it.”—That is the point. Therefore a thankful person is one who sees and knows.—Who never forgot to give thanks? The Lord Jesus.—Prove this from the Catechism. “Our Lord Jesus Christ . . . took bread, gave thanks, brake it,” etc.—So we have another Chief Part that mentions the giving of thanks! Which one is it? The Fifth.—Did Jesus thank only for the bread? No; “after the same manner also He took the cup, when He had supped, gave thanks,” etc.—For what did Jesus give thanks? For everything.—But to thank with words or with the tongue is not all of thanksgiving. How should our gratitude express itself? Think of the First Article. For all of which I am in duty bound to thank and praise, to serve and obey Him.—Service! That is what you owe also to the Savior. For what purpose has He purchased you? “That I might be His own and live under Him in His kingdom and serve Him.”—You also owe this to men. The two parts, “serve and obey,” stand beside each other elsewhere in the Catechism. Where? In the Fourth Commandment. “Serve and obey them, hold them in love and esteem.”—Should you obey everybody? No; only parents and masters. But “serve”? Does that only apply to parents? No; to all men. We should also be of service to our neighbor. Where are we told this? In the Ninth Commandment: “But help and serve him in keeping it.”—To what does this here refer? To his inheritance or house.—Which other Commandment emphasizes this duty? The Seventh: “But help him to improve and protect his property and living.”—When especially are you bound to help him? In every need.—Where is that written? In the Fifth Commandment: “but help and befriend,” etc.—But what if the person be your enemy? “So we on our part will heartily forgive and readily do good to those who sin against us.”—Now we have learned something about thanking and serving from four Chief Parts. Which is lacking? The Fourth Chief Part—of Holy Baptism.—You will not find the word “thank” there, but the idea and thought is there of the very best thanksgiving, which is at the same time a daily serving. Which is that? “That the old Adam in us by . . . . . . . . . putrify forever.”

\quellseite{633}
\section{3. Instruction in Hymns}

The hymn: “\emph{Wake, arise! The voice is calling.}” (According to Schueren, page 119 ff.)

Twice this week the bell in the tower was tolled. Our fathers experienced times when the death knell was sounded every day in the year. Every day the bell tolled its message of death. In the century in which Luther lived, three times a fearful pestilence stalked through the German land—a pestilence so dreadful that in one visitation more than half of the inhabitants in many places died. Those were the unfortunate years 1529, 1572, and 1597. When the last pestilence swept through, there lived in Unna, a small city in the province of Westphalia, a pious Lutheran preacher named Nicolai. He had already suffered much in many places at the hands of the enemies of the followers of Luther. Then came the great pest. Daily he saw from 20 to 30 corpses carried past his house. He spent his time from early morning until late night among the dead and dying. In a short time 1400 persons were buried. Of what could he think but of death and the grave? It was midnight also for his congregation. The Bridegroom was near at hand to lead the wise virgins to the heavenly courts. All things indicated this. He himself was a watchman placed by God upon the towers to remind the congregation of the Bridegroom’s coming. And he was faithful to his trust. He admonished to repentance and then comforted the suffering with the prospect of the heavenly Jerusalem, the city of God with the pearly gates, and put it all in the form of a beautiful hymn. (Teacher recites the hymn.) Now open your books. Sing the first verse. How beautiful the melody and how appropriate to the words! Let us study the words of this hymn!

I will read the first verse of the hymn once more. A capable child then reads it, and then all in chorus. Of which parable does this verse remind you? It reminds of the Parable of the Ten Virgins.—Repeat this parable!—To what are the virgins in this verse admonished? They are admonished to be ready.—For what should they be ready? To meet the bridegroom.—How must the virgins be in order to meet the bridegroom? They must be awake.—What, then, do the watchmen do who are placed by God upon the towers? They call: Awake!—To whom do they thus call in this verse? They call to the city of Jerusalem.—Who is meant thereby? The people are meant thereby.—We shall see what people are meant. Who is the bridegroom? Christ is the bridegroom.—Whom does He wish to lead to the heavenly palace? Those who believe in Him and follow Him.—How do we call
\quellseite{634}
those people who do this? We call them Christians.—What are we therefore to understand by Jerusalem? The Christian congregation.—What is said to it? To awake and be ready.—Why? Because the bridegroom is coming.—When will He come? He will come at midnight.—When will that be? When may the Lord call you away? He may call me away at any hour.—Which hour, then, is the midnight hour? Every hour is the midnight hour.—When should I, therefore, be ready? I should always be ready.—What is meant by being ready and having one’s lamp filled and burning, you have already learned.

I shall read the first stanza again.—You read it!—And you!—And you!—Now the soprano! The alto! Right; now the second stanza should be easier. Treatment of the two following stanzas as before, but more briefly. In the second stanza attention must be called to the fact that the second part is a joyful response to the first part; and in the middle of the third stanza the teacher reads the corresponding Scripture verse from Revelation 21.—At the close the teacher reads the whole. Then the children read it together. The singing of the third stanza closes the lesson. It will hardly be necessary to tell the children to learn the hymn at home. Many already know it by memory, others will learn it without being commanded. During the review one can tell of the blessings which the hymn has brought into the life of many a person.
